% Luc Destombe --- licence CC BY-NC-SA 4.0
% https://creativecommons.org/licenses/by-nc-sa/4.0/deed.fr
% Fichier autonome : compiler avec pdflatex, deux fois (signets, nombre de pages).
\documentclass[11pt,a4paper]{article}
\makeatletter
\newif\iflycee@unecolonne
\lycee@unecolonnetrue
\newif\iflycee@palatino
\lycee@palatinotrue

\RequirePackage[utf8]{inputenc}
\RequirePackage[T1]{fontenc}
\RequirePackage[french]{babel}
\RequirePackage{etoolbox}

\newcommand{\ShorthandsOff}{\@ifundefined{shorthandoff}{}{\shorthandoff{;:!?}}}
\newcommand{\ShorthandsOn}{\@ifundefined{shorthandon}{}{\shorthandon{;:!?}}}
\ShorthandsOff
\AtBeginDocument{\ShorthandsOn}
\AtBeginEnvironment{tikzpicture}{\ShorthandsOff}

\RequirePackage{geometry}
\RequirePackage{fancyhdr}
\RequirePackage{lastpage}
\RequirePackage{multicol}
\RequirePackage{array}
\RequirePackage{enumitem}
\RequirePackage{xcolor}
\RequirePackage{needspace}

\RequirePackage{amsmath,amssymb}
\RequirePackage{mathpazo}
\DeclareMathSizes{10.95}{10.95}{8}{5.5}
\begingroup\catcode`\_=\active \catcode`\^=\active
\gdef_#1{\sb{\scriptscriptstyle #1}}
\gdef^#1{\sp{\scriptscriptstyle\lycee@moinsepais #1}}
\endgroup
\newcommand\lycee@moins{\mathbin{%
  \setbox\z@\hbox{$\m@th\scriptscriptstyle\lycee@moinsorig$}%
  \dimen@=.55\wd\z@ \advance\dimen@-.5pt
  \hbox to.75\wd\z@{\hss$\m@th\scriptscriptstyle\vcenter{\hbox{%
    \tikz\draw[line width=.5pt,line cap=round](0,0)--(\the\dimen@,0);}}$\hss}}}
\begingroup\lccode`\~=`\- \lowercase{\endgroup\let~}\lycee@moins
\newcommand\lycee@moinsepais{\mathcode`\-="8000 }
\AtBeginDocument{\mathchardef\lycee@moinsorig=\mathcode`\-\relax}
\AtBeginDocument{\catcode`\_=12 \mathcode`\_="8000
  \catcode`\^=12 \mathcode`\^="8000 }
\newcommand\lycee@exposants{%
  \fontdimen13\textfont2=.48\fontdimen6\textfont2
  \fontdimen14\textfont2=.48\fontdimen6\textfont2
  \fontdimen15\textfont2=.48\fontdimen6\textfont2
  \fontdimen13\scriptfont2=.48\fontdimen6\scriptfont2
  \fontdimen14\scriptfont2=.48\fontdimen6\scriptfont2
  \fontdimen15\scriptfont2=.48\fontdimen6\scriptfont2}
\every@math@size\expandafter{\the\every@math@size\lycee@exposants}
\newlength{\retraitcalculs}
\setlength{\retraitcalculs}{2em}

\RequirePackage{tikz}
\usetikzlibrary{arrows.meta,calc,babel,shapes.misc}
\RequirePackage[most]{tcolorbox}
\tcbuselibrary{skins,breakable}

\definecolor{coulDef}{HTML}{1F4E79}
\definecolor{coulProp}{HTML}{7030A0}
\definecolor{coulRem}{HTML}{C55A11}
\definecolor{coulEx}{HTML}{2E7D32}
\definecolor{coulExo}{HTML}{B03A2E}
\definecolor{coulPart}{HTML}{1F4E79}
\definecolor{coulMeth}{HTML}{1B6E4F}
\definecolor{coulCourbe}{HTML}{0B5ED7}
\definecolor{coulCourbeB}{HTML}{C00000}
\definecolor{coulCourbeC}{HTML}{2E7D32}
\definecolor{coulGrille}{HTML}{8C8C8C}
\definecolor{coulRep}{HTML}{1B6E4F}
\definecolor{coulBar}{HTML}{2E7D32}

\newcommand{\R}{\mathbb{R}}
\newcommand{\Rs}{\mathbb{R}^{*}}
\renewcommand{\le}{\leqslant}
\renewcommand{\ge}{\geqslant}
\renewcommand{\approx}{\simeq}
\newcommand{\vect}[1]{\overrightarrow{#1}}
\newcommand{\norme}[1]{\left\lVert #1 \right\rVert}
\newcommand{\intervalle}[2]{\left[#1\,;#2\right]}
\RequirePackage{cancel}
\renewcommand{\CancelColor}{\color{coulExo}}
\newcommand{\fois}[1]{\times{\color{coulExo}#1}}
\newcommand{\barre}[1]{\cancel{#1}}

\tikzset{grille/.style={coulGrille,line width=0.4pt}}
\newcommand{\quadrillage}[4]{\draw[grille,step=1] (#1,#2) grid (#3,#4);}
\tikzset{
  croix/.style={cross out,draw,line width=0.6pt,inner sep=0pt,minimum size=4pt},
  croixdroite/.style={croix,rotate=45,line width=0.8pt},
}
\newcommand{\pt}[3]{\path (#1) node[croix]{} node[#2,font=\small,inner sep=2.5pt]{$#3$};}

\newcommand{\lycee@repere}[4]{%
  \draw[grille,step=1] (#1,#2) grid (#3,#4);
  \draw[-{Stealth[length=2mm]},thick] (#1,0)--({#3+0.4},0) node[below left=-1pt]{$x$};
  \draw[-{Stealth[length=2mm]},thick] (0,#2)--(0,{#4+0.4}) node[below left=-1pt]{$y$};
  \node[below left,font=\scriptsize,inner sep=1.5pt] at (0,0) {$0$};}
\newcommand{\lycee@gradx}[1]{\foreach \k/\l in {#1}{%
  \draw (\k,0.07)--(\k,-0.07) node[below,font=\scriptsize,inner sep=1.5pt]{$\l$};}}
\newcommand{\lycee@grady}[1]{\foreach \k/\l in {#1}{%
  \draw (0.07,\k)--(-0.07,\k) node[left,font=\scriptsize,inner sep=1.5pt]{$\l$};}}
\newcommand{\lycee@intff}[2]{\left[#1\,;#2\right]}
\newcommand{\lycee@intfo}[2]{\left[#1\,;#2\right[}
\newcommand{\lycee@intof}[2]{\left]#1\,;#2\right]}
\newcommand{\lycee@intoo}[2]{\left]#1\,;#2\right[}
\newcommand{\lycee@ens}[1]{\left\{#1\right\}}
\AtBeginDocument{%
  \providecommand{\repere}{\lycee@repere}%
  \providecommand{\gradx}{\lycee@gradx}%
  \providecommand{\grady}{\lycee@grady}%
  \providecommand{\Cf}{\mathcal{C}\sb{\scriptscriptstyle f}}%
  \providecommand{\Cg}{\mathcal{C}\sb{\scriptscriptstyle g}}%
  \providecommand{\intff}{\lycee@intff}%
  \providecommand{\intfo}{\lycee@intfo}%
  \providecommand{\intof}{\lycee@intof}%
  \providecommand{\intoo}{\lycee@intoo}%
  \providecommand{\ens}{\lycee@ens}}

\newlist{questions}{enumerate}{3}
\setlist[questions,1]{label=\arabic*\textdegree),leftmargin=*,itemsep=2pt,topsep=3pt}
\setlist[questions,2]{label=\alph*),leftmargin=*,itemsep=1pt,topsep=2pt}
\setlist[questions,3]{label=\roman*),leftmargin=*,itemsep=1pt,topsep=1pt}
\setlist[itemize]{label=\textbullet,leftmargin=*,itemsep=2pt,topsep=3pt}

\newcommand{\besoinplace}[1]{\par\penalty\@M\begingroup
  \dimen@=#1\relax\dimen@ii=\pagegoal\advance\dimen@ii-\pagetotal
  \ifdim\dimen@>\dimen@ii\ifdim\dimen@ii>\z@\vfil\fi\break\fi\endgroup}

\newcommand{\lycee@pied}{\thepage/\pageref{LastPage}}
\newcommand{\entete}[2]{%
  \pagestyle{fancy}\fancyhf{}%
  \fancyhead[L]{\footnotesize\color{coulPart}\textbf{#1}}%
  \fancyhead[R]{\footnotesize\color{coulPart}\textbf{#2}}%
  \fancyfoot[C]{\footnotesize\lycee@pied}%
  \renewcommand{\headrulewidth}{0.4pt}%
  \renewcommand{\headrule}{\hbox to\headwidth{%
    \color{coulPart!40}\leaders\hrule height \headrulewidth\hfill}}}

\RequirePackage{ccicons}
\newcommand{\lycee@licence}{{\scriptsize\color{gray}%
  \copyright\ Luc Destombe --- \ccbyncsa\ CC BY-NC-SA 4.0}}
\newif\iflycee@licence \lycee@licencetrue
\newcommand{\sanslicence}{\lycee@licencefalse}
\AtBeginDocument{\iflycee@licence\AddToHookNext{shipout/foreground}{%
  \put(0,0){\rlap{\kern\dimexpr1in+\hoffset+\oddsidemargin+\textwidth\relax
    \llap{\raisebox{-\dimexpr1in+\voffset+\topmargin+\headheight+\headsep
      +\textheight+\footskip\relax}{\lycee@licence}}}}}\fi}

\newcommand{\formulesserrees}{%
  \setlength{\abovedisplayskip}{3pt}\setlength{\belowdisplayskip}{3pt}%
  \setlength{\abovedisplayshortskip}{2pt}\setlength{\belowdisplayshortskip}{3pt}}

\setlength{\parindent}{0pt}

\RequirePackage[implicit=false,hidelinks,pdfencoding=auto,psdextra,bookmarksopen,
  bookmarksopenlevel=1,pdfcreator={},pdfproducer={}]{hyperref}
\RequirePackage{bookmark}
\hypersetup{pdfauthor={Luc Destombe},pdfkeywords={CC BY-NC-SA 4.0}}
\newcounter{lycee@signet}
\newcommand{\signet}[3][]{\stepcounter{lycee@signet}%
  \edef\lycee@dest{\if\relax\detokenize{#1}\relax
    signet.\arabic{lycee@signet}\else#1\fi}%
  \hypertarget{\lycee@dest}{}\bookmark[level=#2,dest=\lycee@dest]{#3}}
\pdfstringdefDisableCommands{%
  \def\R{R}\def\vect#1{#1}\def\Cf{Cf}\def\Cg{Cg}\def\fois#1{×#1}%
  \def\barre#1{#1}\def\intervalle#1#2{[#1;#2]}\def\intff#1#2{[#1;#2]}%
  \def\textsuperscript#1{#1}\def\dfrac#1#2{#1/#2}\def\frac#1#2{#1/#2}%
  \def\vec#1{#1⃗}}%
\RequirePackage[official]{eurosym}

\tcbset{exercice corrige/.style={
  enhanced,breakable,before upper=\raggedright,
  colback=white,colframe=coulExo,
  boxrule=0.8pt,arc=2pt,
  left=6pt,right=6pt,top=8pt,bottom=5pt,
  before skip=9pt,after skip=4pt,
  coltitle=white,fonttitle=\bfseries\small,
  attach boxed title to top left={xshift=8pt,yshift=-\tcboxedtitleheight/2},
  boxed title style={colback=coulExo,arc=2pt,boxrule=0pt}}}
\newcounter{exo}
\newtcolorbox{exercice}[1][]{exercice corrige,
  phantom={\signet[exercice-\arabic{exo}]{2}{Exercice \arabic{exo}}},
  title={Exercice \theexo\ifx\relax#1\relax\else{} --- #1\fi}}
\newenvironment{exo}[1][\relax]{\refstepcounter{exo}\begin{exercice}[#1]}{\end{exercice}}

\geometry{top=1.8cm,bottom=1cm,left=1cm,right=1cm,headheight=14pt,footskip=0.6cm}
\renewcommand{\lycee@pied}{\thepage}

\renewcommand{\theexo}{\ifnum\value{exo}<10 0\fi\arabic{exo}}
\tcbset{exercice corrige/.style={
  enhanced,breakable,before upper=\raggedright,
  colback=white,colframe=coulExo,
  boxrule=0.9pt,arc=2pt,
  left=8pt,right=8pt,top=8pt,bottom=6pt,
  before skip=8pt,after skip=6pt,
  coltitle=white,fonttitle=\bfseries,
  attach boxed title to top left={xshift=8pt,yshift=-\tcboxedtitleheight/2},
  boxed title style={colback=coulExo,arc=2pt,boxrule=0pt}}}

\newcommand{\entetefiche}[2]{%
  \begin{center}
    \begin{tcolorbox}[enhanced,width=0.92\linewidth,colback=white,colframe=coulPart,
        boxrule=1.4pt,arc=3pt,halign=center,top=6pt,bottom=6pt]
      {\large\bfseries\color{coulPart} CORRIGÉ --- FICHE D'EXERCICES}\\[3pt]
      {\normalsize\bfseries #1}\\[2pt]
      {\normalsize #2}
    \end{tcolorbox}
  \end{center}}

\newcounter{partie}
\newcommand{\partie}[1]{%
  \stepcounter{partie}%
  \besoinplace{8\baselineskip}\vspace{4pt}%
  \begin{tcolorbox}[enhanced,colback=coulPart!8,colframe=coulPart,boxrule=0pt,
      leftrule=3pt,arc=0pt,left=6pt,right=4pt,top=3pt,bottom=3pt,
      before skip=6pt,after skip=2pt]
    \bfseries\color{coulPart}\Roman{partie} --- #1\signet{1}{\Roman{partie} --- #1}
  \end{tcolorbox}}
\newcommand{\partiesansnum}[1]{%
  \besoinplace{8\baselineskip}\vspace{4pt}%
  \begin{tcolorbox}[enhanced,colback=coulPart!8,colframe=coulPart,boxrule=0pt,
      leftrule=3pt,arc=0pt,left=6pt,right=4pt,top=3pt,bottom=3pt,
      before skip=6pt,after skip=2pt]
    \bfseries\color{coulPart}#1\signet{1}{#1}
  \end{tcolorbox}}

\setlist[questions,1]{label=\arabic*\textdegree),leftmargin=*,itemsep=3pt,topsep=3pt}
\setlist[questions,2]{label=\alph*),leftmargin=*,itemsep=2pt,topsep=2pt}
\setlist[questions,3]{label=\roman*),leftmargin=*,itemsep=1pt,topsep=1pt}
\newlist{expr}{enumerate}{1}
\setlist[expr]{label={\Alph*\ $=$},leftmargin=*,itemsep=2pt,topsep=3pt}

\newcommand{\res}[1]{{\color{coulRep}#1}}
\newcommand{\rem}[1]{\par\smallskip{\small\itshape\color{black!60}$\rightarrow$ #1}}
\newenvironment{calculs}{\par\smallskip\noindent$\displaystyle\begin{aligned}[t]}%
  {\end{aligned}$\par\smallskip}
\newcommand{\justif}[1]{\quad\text{\small\itshape\color{black!60}#1}}
\newcommand{\cas}[1]{\par\smallskip\noindent\textbf{\color{coulExo}#1}\nobreak}
\newcommand{\explique}[2]{\par\smallskip\noindent
  \begin{minipage}[t]{0.6\linewidth}\raggedright #1\end{minipage}\hfill
  \begin{minipage}[t]{0.37\linewidth}\raggedright\leavevmode
    {\small\itshape\color{black!60}#2}\end{minipage}\par}

\setlength{\parskip}{2pt}

\makeatother
\usepackage{tkz-tab}
\entete{Chapitre 5 --- Dérivation et étude de fonctions --- corrigé}{1\textsuperscript{re} STI2D}
\let\casligne\cas
\renewcommand{\cas}[1]{\casligne{#1}\par\nopagebreak}
\setlength{\parindent}{0pt}
\tikzset{lecture/.style={coulRep,line width=1pt,-{Stealth[length=2mm]}}}

\begin{document}

\entetefiche{Chapitre 5 --- Dérivation et étude de fonctions}{Première STI2D}

\partie{Fonctions affines et droites}

\begin{exo}[Coefficient directeur]
\explique{Dans chaque cas, on recopie la formule du coefficient directeur, puis on
remplace par les coordonnées.}{ce qu'on monte, divisé par ce qu'on avance}
\cas{a)}
\begin{calculs}
m &= \frac{y_{B}-y_{A}}{x_{B}-x_{A}} &&\justif{on écrit la formule}\\
m &= \frac{8-2}{3-1} &&\justif{on remplace par les coordonnées}\\
m &= \frac{6}{2} &&\\
m &= \res{3} &&
\end{calculs}
\cas{b)}
\begin{calculs}
m &= \frac{-2-4}{2-(-1)} &&\justif{$2-(-1)=2+1$ : attention aux signes}\\
m &= \frac{-6}{3} &&\\
m &= \res{-2} &&\justif{la droite descend}
\end{calculs}
\cas{c)}
\begin{calculs}
m &= \frac{-1-(-3)}{4-0} &&\justif{$-1-(-3)=-1+3$}\\
m &= \frac{2}{4} &&\\
m &= \res{\frac{1}{2}} &&\justif{on simplifie par $2$}
\end{calculs}
\cas{d)}
\begin{calculs}
m &= \frac{5-5}{1-(-3)} &&\\
m &= \frac{0}{4} &&\\
m &= \res{0} &&
\end{calculs}
\explique{$A$ et $B$ ont la même ordonnée : la droite $(AB)$ est horizontale.}{on ne
monte pas}
\rem{Erreur fréquente : soustraire dans un ordre en haut et dans l'autre en bas. Si l'on
commence par $B$ au numérateur, on commence aussi par $B$ au dénominateur.}
\end{exo}

\begin{exo}[Équation d'une droite]
\explique{Comme à la méthode 1 : on calcule le coefficient directeur $m$, puis on trouve
$p$ avec le point $A$.}{méthode 1}
\cas{a)}
\begin{calculs}
m &= \frac{y_{B}-y_{A}}{x_{B}-x_{A}} &&\justif{on écrit la formule}\\
m &= \frac{4-1}{2-1} &&\justif{on remplace par les coordonnées}\\
m &= 3 &&
\end{calculs}
\explique{L'équation est de la forme $y=3x+p$. La droite passe par $A(1\,;1)$ :}{on
cherche maintenant $p$}\nopagebreak
\begin{calculs}
3\times 1+p &= 1 &&\justif{on remplace $x$ par $1$ et $y$ par $1$}\\
p &= 1-3 &&\\
p &= -2 &&
\end{calculs}
\explique{La droite $(AB)$ a pour équation $\res{y=3x-2}$.}{vérification avec $B$ :
$3\times 2-2=4$}
\cas{b)}
\begin{calculs}
m &= \frac{1-5}{1-(-1)} &&\justif{attention aux signes}\\
m &= \frac{-4}{2} &&\\
m &= -2 &&
\end{calculs}
\explique{L'équation est de la forme $y=-2x+p$. La droite passe par $A(-1\,;5)$ :}{}\nopagebreak
\begin{calculs}
-2\times(-1)+p &= 5 &&\justif{on remplace $x$ par $-1$ et $y$ par $5$}\\
2+p &= 5 &&\\
p &= 3 &&
\end{calculs}
\explique{La droite $(AB)$ a pour équation $\res{y=-2x+3}$.}{vérification avec $B$ :
$-2\times 1+3=1$}
\cas{c)}
\begin{calculs}
m &= \frac{5-(-3)}{2-(-2)} &&\justif{$5+3$ en haut, $2+2$ en bas}\\
m &= \frac{8}{4} &&\\
m &= 2 &&
\end{calculs}
\explique{L'équation est de la forme $y=2x+p$. La droite passe par $A(-2\,;-3)$ :}{}\nopagebreak
\begin{calculs}
2\times(-2)+p &= -3 &&\justif{on remplace $x$ par $-2$ et $y$ par $-3$}\\
-4+p &= -3 &&\\
p &= 1 &&
\end{calculs}
\explique{La droite $(AB)$ a pour équation $\res{y=2x+1}$.}{vérification avec $B$ :
$2\times 2+1=5$}
\cas{d)}
\begin{calculs}
m &= \frac{2-0}{6-2} &&\\
m &= \frac{2}{4} &&\\
m &= \frac{1}{2} &&\justif{on simplifie par $2$}
\end{calculs}
\explique{L'équation est de la forme $y=\frac{1}{2}x+p$. La droite passe par
$A(2\,;0)$ :}{}\nopagebreak
\begin{calculs}
\frac{1}{2}\times 2+p &= 0 &&\justif{on remplace $x$ par $2$ et $y$ par $0$}\\
1+p &= 0 &&\justif{$\frac{1}{2}\times 2=1$}\\
p &= -1 &&
\end{calculs}
\explique{La droite $(AB)$ a pour équation $\res{y=\frac{1}{2}x-1}$.}{vérification avec
$B$ : $\frac{1}{2}\times 6-1=2$}
\rem{Le second point sert de contrôle : ses coordonnées doivent vérifier l'équation
trouvée.}
\end{exo}

\begin{exo}[Une fonction affine par deux images]
\explique{$f(2)=7$ veut dire que la droite qui représente $f$ passe par le point
$(2\,;7)$ : on se ramène à l'exercice 2.}{une image, c'est un point de la droite}
\cas{1)}
\explique{La droite passe par $(2\,;7)$ et $(-1\,;-2)$.}{$f(2)=7$ et $f(-1)=-2$}\nopagebreak
\begin{calculs}
m &= \frac{7-(-2)}{2-(-1)} &&\justif{la formule du coefficient directeur}\\
m &= \frac{9}{3} &&\\
m &= 3 &&
\end{calculs}
\explique{Donc $f(x)=3x+p$, avec $f(2)=7$ :}{on cherche $p$}\nopagebreak
\begin{calculs}
3\times 2+p &= 7 &&\justif{on remplace $x$ par $2$}\\
6+p &= 7 &&\\
p &= 1 &&
\end{calculs}
\explique{Donc $\res{f(x)=3x+1}$.}{vérification : $f(-1)=-3+1=-2$}\nopagebreak
\begin{calculs}
f(10) &= 3\times 10+1 &&\justif{on remplace $x$ par $10$}\\
f(10) &= \res{31} &&
\end{calculs}
\cas{2)}
\explique{La droite passe par $(-3\,;8)$ et $(1\,;0)$.}{$g(-3)=8$ et $g(1)=0$}\nopagebreak
\begin{calculs}
m &= \frac{0-8}{1-(-3)} &&\\
m &= \frac{-8}{4} &&\\
m &= -2 &&
\end{calculs}
\explique{Donc $g(x)=-2x+p$, avec $g(1)=0$ :}{}\nopagebreak
\begin{calculs}
-2\times 1+p &= 0 &&\justif{on remplace $x$ par $1$}\\
p &= 2 &&
\end{calculs}
\explique{Donc $\res{g(x)=-2x+2}$.}{vérification : $g(-3)=6+2=8$}\nopagebreak
\begin{calculs}
g(5) &= -2\times 5+2 &&\\
g(5) &= \res{-8} &&
\end{calculs}
\end{exo}

\begin{exo}[Lire des équations]
\begin{minipage}[c]{0.56\linewidth}\raggedright
Pour chaque droite, on lit $p$ là où elle coupe l'axe des ordonnées (croix), puis on
part de ce point : on avance et on regarde de combien on monte (flèches vertes).

\smallskip
$(d_{1})$ : elle coupe l'axe en $-1$, donc $p=-1$ ; on avance de $1$, on monte de $2$,
donc $m=2$ : $\res{y=2x-1}$.

\smallskip
$(d_{2})$ : $p=3$ ; on avance de $1$, on descend de $1$, donc $m=-1$ :
$\res{y=-x+3}$.

\smallskip
$(d_{3})$ : $p=-2$ ; on avance de $2$, on monte de $1$, donc $m=\frac{1}{2}$ :
$\res{y=\frac{1}{2}x-2}$.
\end{minipage}\hfill
\begin{minipage}[c]{0.4\linewidth}\centering
\begin{tikzpicture}[x=0.65cm,y=0.65cm]
  \repere{-3}{-3}{5}{4}
  \gradx{-2,-1,1,2,3,4} \grady{-2,-1,1,2,3}
  \begin{scope}
  \clip (-3,-3) rectangle (5,4);
  \draw[coulCourbe,line width=1.1pt] (-1,-3)--(2.5,4);
  \draw[coulCourbeB,line width=1.1pt] (-1,4)--(5,-2);
  \draw[coulCourbeC,line width=1.1pt] (-2,-3)--(5,0.5);
  \end{scope}
  \draw[lecture] (0,-1)--(1,-1);
  \draw[lecture] (1,-1)--(1,1);
  \draw[lecture] (0,3)--(1,3);
  \draw[lecture] (1,3)--(1,2);
  \draw[lecture] (0,-2)--(2,-2);
  \draw[lecture] (2,-2)--(2,-1);
  \path[coulCourbe] (0,-1) node[croix]{};
  \path[coulCourbeB] (0,3) node[croix]{};
  \path[coulCourbeC] (0,-2) node[croix]{};
  \node[coulCourbe,font=\footnotesize] at (1.4,3.6) {$(d_{1})$};
  \node[coulCourbeB,font=\footnotesize] at (-0.6,2.7) {$(d_{2})$};
  \node[coulCourbeC,font=\footnotesize] at (4.2,0.8) {$(d_{3})$};
\end{tikzpicture}
\end{minipage}
\rem{Pour $(d_{3})$, avancer de $1$ ne tombe pas sur un nœud du quadrillage : on avance
de $2$ et on divise ce qu'on monte par $2$.}
\end{exo}

\partie{Nombre dérivé}

\begin{exo}[Nombre dérivé en $1$]
\cas{a)}
\begin{calculs}
f(1) &= 1^{2}+2\times 1-1 &&\justif{on remplace $x$ par $1$}\\
f(1) &= \res{2} &&\\[4pt]
f(1+h) &= (1+h)^{2}+2(1+h)-1 &&\justif{on remplace $x$ par $(1+h)$}\\
f(1+h) &= 1+2h+h^{2}+2+2h-1 &&\justif{$(a+b)^{2}=a^{2}+2ab+b^{2}$, puis on développe}\\
f(1+h) &= \res{h^{2}+4h+2} &&\justif{on réduit}
\end{calculs}
\cas{b)}
\begin{calculs}
T(h) &= \frac{f(1+h)-f(1)}{h} &&\justif{la formule, avec $a=1$}\\
T(h) &= \frac{h^{2}+4h+2-2}{h} &&\justif{on remplace}\\
T(h) &= \frac{h^{2}+4h}{h} &&\justif{$2-2=0$}\\
T(h) &= \frac{\barre{h}\,(h+4)}{\barre{h}} &&\justif{on factorise par $h$, puis on simplifie}\\
T(h) &= \res{h+4} &&
\end{calculs}
\cas{c)}
\explique{Quand $h$ se rapproche de $0$, $T(h)=h+4$ se rapproche de $4$. Donc
$\res{f'(1)=4}$.}{on remplace $h$ par $0$ dans le résultat simplifié}
\rem{Contrôle : au numérateur de $T(h)$, les nombres seuls s'annulent toujours ($2-2$) ;
s'il en reste un, il y a une erreur de calcul.}
\end{exo}

\begin{exo}[Nombre dérivé en $2$]
\cas{Résolution}
\explique{Les deux images d'abord :}{méthode 2}\nopagebreak
\begin{calculs}
g(2) &= 2\times 2^{2}-3\times 2 &&\justif{on remplace $x$ par $2$}\\
g(2) &= 8-6 &&\\
g(2) &= \res{2} &&\\[4pt]
g(2+h) &= 2(2+h)^{2}-3(2+h) &&\justif{on remplace $x$ par $(2+h)$}\\
g(2+h) &= 2(4+4h+h^{2})-6-3h &&\justif{on développe d'abord le carré}\\
g(2+h) &= 8+8h+2h^{2}-6-3h &&\justif{puis on multiplie par $2$}\\
g(2+h) &= \res{2h^{2}+5h+2} &&\justif{on réduit}
\end{calculs}
\explique{Puis le taux de variation :}{}\nopagebreak
\begin{calculs}
T(h) &= \frac{g(2+h)-g(2)}{h} &&\justif{la formule, avec $a=2$}\\
T(h) &= \frac{2h^{2}+5h+2-2}{h} &&\justif{on remplace}\\
T(h) &= \frac{\barre{h}\,(2h+5)}{\barre{h}} &&\justif{on factorise par $h$, puis on simplifie}\\
T(h) &= \res{2h+5} &&
\end{calculs}
\explique{Quand $h$ se rapproche de $0$, $T(h)=2h+5$ se rapproche de $5$. Donc
$\res{g'(2)=5}$.}{$2h$ se rapproche de $0$}
\rem{Erreur fréquente : $2(2+h)^{2}=(4+2h)^{2}$. Le carré passe avant la multiplication
par $2$ (priorité des puissances).}
\end{exo}

\begin{exo}[Avec un coefficient négatif]
\cas{Résolution}
\begin{calculs}
k(1) &= -1^{2}+5\times 1 &&\justif{$-1^{2}=-1$ : le carré porte sur $1$ seulement}\\
k(1) &= \res{4} &&\\[4pt]
k(1+h) &= -(1+h)^{2}+5(1+h) &&\justif{on remplace $x$ par $(1+h)$}\\
k(1+h) &= -(1+2h+h^{2})+5+5h &&\justif{le $-$ porte sur tout le carré}\\
k(1+h) &= -1-2h-h^{2}+5+5h &&\justif{on enlève les parenthèses : les signes changent}\\
k(1+h) &= \res{-h^{2}+3h+4} &&\justif{on réduit}
\end{calculs}
\begin{calculs}
T(h) &= \frac{k(1+h)-k(1)}{h} &&\justif{la formule, avec $a=1$}\\
T(h) &= \frac{-h^{2}+3h+4-4}{h} &&\justif{on remplace}\\
T(h) &= \frac{\barre{h}\,(-h+3)}{\barre{h}} &&\justif{on factorise par $h$, puis on simplifie}\\
T(h) &= \res{-h+3} &&
\end{calculs}
\explique{Quand $h$ se rapproche de $0$, $T(h)=-h+3$ se rapproche de $3$. Donc
$\res{k'(1)=3}$.}{on remplace $h$ par $0$}
\rem{Erreur fréquente : $-(1+h)^{2}=-1+2h+h^{2}$. Le signe $-$ change les trois termes du
développement.}
\end{exo}

\begin{exo}[À la calculatrice]
\cas{a)}
\explique{Pour chaque valeur de $h$, on calcule $f(4+h)=(4+h)^{2}$, on retire
$f(4)=16$, puis on divise par $h$. Détail pour $h=0{,}1$ :}{on peut taper directement
l'expression de $T(h)$}\nopagebreak
\begin{calculs}
T(0{,}1) &= \frac{4{,}1^{2}-16}{0{,}1} &&\justif{$4+0{,}1=4{,}1$}\\
T(0{,}1) &= \frac{16{,}81-16}{0{,}1} &&\\
T(0{,}1) &= \frac{0{,}81}{0{,}1} &&\\
T(0{,}1) &= 8{,}1 &&\justif{diviser par $0{,}1$, c'est multiplier par $10$}
\end{calculs}
\explique{De même pour les deux autres valeurs :}{}\nopagebreak
\begin{calculs}
T(1) &= \frac{25-16}{1} = 9 &&\justif{$5^{2}=25$}\\
T(0{,}01) &= \frac{16{,}0801-16}{0{,}01} = 8{,}01 &&\justif{$4{,}01^{2}=16{,}0801$}
\end{calculs}
\begin{center}
\renewcommand{\arraystretch}{1.4}
\begin{tabular}{|c|c|c|c|}
\hline
$h$ & $1$ & $0{,}1$ & $0{,}01$\\
\hline
$T(h)$ & $\res{9}$ & $\res{8{,}1}$ & $\res{8{,}01}$\\
\hline
\end{tabular}
\end{center}
\cas{b)}
\explique{Quand $h$ se rapproche de $0$, $T(h)$ se rapproche de $8$ : le tableau
suggère $\res{f'(4)=8}$.}{$9$, puis $8{,}1$, puis $8{,}01$ : l'écart à $8$ est divisé
par $10$ à chaque fois}
\end{exo}

\partie{Tangente à une courbe}

\begin{exo}[Lire des nombres dérivés]
\begin{center}
\begin{tikzpicture}[x=0.65cm,y=0.65cm]
  \repere{-4}{-2}{5}{4}
  \gradx{-3,-2,-1,1,2,3,4} \grady{-1,1,2,3}
  \draw[coulCourbe,line width=1.1pt]
    (-3,2) .. controls (-2.333,0.667) and (-1.667,-1) .. (-1,-1)
           .. controls (-0.333,-1) and (0.333,-0.333) .. (1,1)
           .. controls (1.667,2.333) and (2.333,3) .. (3,3)
           .. controls (3.333,3) and (3.667,2.667) .. (4,2);
  \draw[coulCourbeB,thick] (-4,4)--(-1.7,-0.6);
  \draw[coulCourbeB,thick] (-2,-1)--(0,-1);
  \draw[coulCourbeB,thick] (0,-1)--(2.5,4);
  \draw[coulCourbeB,thick] (3,4)--(4.9,0.2);
  \draw[lecture] (-3,2)--(-2,2);
  \draw[lecture] (-2,2)--(-2,0);
  \draw[lecture] (1,1)--(2,1);
  \draw[lecture] (2,1)--(2,3);
  \draw[lecture] (3,4)--(4,4);
  \draw[lecture] (4,4)--(4,2);
  \pt{-3,2}{left}{A}
  \pt{-1,-1}{below}{B}
  \pt{1,1}{left}{C}
  \pt{4,2}{right}{E}
  \node[coulCourbe,font=\footnotesize] at (2.4,2.4) {$\Cf$};
\end{tikzpicture}
\end{center}
\explique{En $A(-3\,;2)$ : on avance de $1$, on descend de $2$. Donc
$f'(-3)=\res{-2}$.}{on descend : nombre négatif}
\explique{En $B(-1\,;-1)$ : la tangente est horizontale. Donc $f'(-1)=\res{0}$.}{on ne
monte pas}
\explique{En $C(1\,;1)$ : on avance de $1$, on monte de $2$. Donc
$f'(1)=\res{2}$.}{la tangente passe par $(2\,;3)$}
\explique{En $E(4\,;2)$ : la tangente passe par $(3\,;4)$ ; de là, on avance de $1$ et
on descend de $2$ jusqu'à $E$. Donc $f'(4)=\res{-2}$.}{on part d'un nœud de la
tangente}
\rem{On lit $f'(a)$ sur la tangente (une pente), pas sur la courbe : en $C$, $f(1)=1$
mais $f'(1)=2$.}
\end{exo}

\begin{exo}[Tracer des tangentes]
\explique{Dans chaque cas, on place le point, puis on part de ce point en suivant le
nombre dérivé : on avance de $1$ et on monte de $f'(a)$. On trace la droite qui passe par
les deux points obtenus.}{méthode 3 ; une figure par tangente}
\cas{a)}
\begin{minipage}[c]{0.58\linewidth}\raggedright
\explique{On place $A(0\,;1)$. On avance de $1$ et on monte de $2$ : on arrive en
$(1\,;3)$.}{$f'(0)=2$}
\explique{On trace la droite qui passe par ces deux points.}{elle passe aussi par
$(-1\,;-1)$}
\end{minipage}\hfill
\begin{minipage}[c]{0.38\linewidth}\centering
\begin{tikzpicture}[x=0.4cm,y=0.4cm]
  \repere{-4}{-4}{4}{4}
  \gradx{-2,2} \grady{-2,2}
  \begin{scope}
  \clip (-4,-4) rectangle (4,4);
  \draw[coulCourbeB,thick] (-2.5,-4)--(1.5,4);
  \end{scope}
  \draw[lecture] (0,1)--(1,1);
  \draw[lecture] (1,1)--(1,3);
  \pt{0,1}{left}{A}
\end{tikzpicture}
\end{minipage}
\cas{b)}
\begin{minipage}[c]{0.58\linewidth}\raggedright
\explique{On place $B(1\,;3)$. On avance de $1$ et on descend de $1$ : on arrive en
$(2\,;2)$.}{$f'(1)=-1$ : on descend}
\explique{On trace la droite qui passe par ces deux points.}{elle passe aussi par
$(4\,;0)$}
\end{minipage}\hfill
\begin{minipage}[c]{0.38\linewidth}\centering
\begin{tikzpicture}[x=0.4cm,y=0.4cm]
  \repere{-4}{-4}{4}{4}
  \gradx{-2,2} \grady{-2,2}
  \begin{scope}
  \clip (-4,-4) rectangle (4,4);
  \draw[coulCourbeB,thick] (-0.5,4.5)--(4.5,-0.5);
  \end{scope}
  \draw[lecture] (1,3)--(2,3);
  \draw[lecture] (2,3)--(2,2);
  \pt{1,3}{left}{B}
\end{tikzpicture}
\end{minipage}
\cas{c)}
\begin{minipage}[c]{0.58\linewidth}\raggedright
\explique{On place $C(-2\,;-1)$. Monter de $\frac{1}{2}$ ne tombe pas sur le
quadrillage : on avance de $2$ et on monte de $1$, on arrive en $(0\,;0)$.}{$f'(-2)=\frac{1}{2}$ :
deux fois plus loin, on monte deux fois plus}
\explique{On trace la droite qui passe par $C$ et par l'origine.}{elle passe aussi par
$(2\,;1)$}
\end{minipage}\hfill
\begin{minipage}[c]{0.38\linewidth}\centering
\begin{tikzpicture}[x=0.4cm,y=0.4cm]
  \repere{-4}{-4}{4}{4}
  \gradx{-2,2} \grady{-2,2}
  \draw[coulCourbeB,thick] (-4,-2)--(4,2);
  \draw[lecture] (-2,-1)--(0,-1);
  \draw[lecture] (0,-1)--(0,0);
  \pt{-2,-1}{below}{C}
\end{tikzpicture}
\end{minipage}
\cas{d)}
\begin{minipage}[c]{0.58\linewidth}\raggedright
\explique{On place $D(3\,;-2)$. Comme $f'(3)=0$, on ne monte pas : la tangente est la
droite horizontale qui passe par $D$.}{pente nulle}
\end{minipage}\hfill
\begin{minipage}[c]{0.38\linewidth}\centering
\begin{tikzpicture}[x=0.4cm,y=0.4cm]
  \repere{-4}{-4}{4}{4}
  \gradx{-2,2} \grady{-2,2}
  \draw[coulCourbeB,thick] (-4,-2)--(4,-2);
  \pt{3,-2}{below}{D}
\end{tikzpicture}
\end{minipage}
\rem{Une tangente se trace avec deux points sur des nœuds du quadrillage : le point de la
courbe, puis le point obtenu en suivant la pente. Un seul point et un tracé \og à
l'\oe il\fg{} ne suffisent pas.}
\end{exo}

\begin{exo}[Retrouver l'allure d'une courbe]
\cas{Résolution}
\begin{center}
\begin{tikzpicture}[x=0.6cm,y=0.6cm]
  \repere{-4}{-4}{4}{3}
  \gradx{-3,-2,-1,1,2,3} \grady{-3,-2,-1,1,2}
  \draw[coulCourbeB,thick] (-3.6,-3.8)--(-2.2,0.4);
  \draw[coulCourbeB,thick] (-2,2)--(0,2);
  \draw[coulCourbeB,thick] (0.3,1.4)--(1.7,-1.4);
  \draw[coulCourbeB,thick] (2,-2)--(4,-2);
  \draw[lecture] (-3,-2)--(-2,-2);
  \draw[lecture] (-2,-2)--(-2,1);
  \draw[lecture] (1,0)--(2,0);
  \draw[lecture] (2,0)--(2,-2);
  \draw[coulCourbe,line width=1.1pt]
    (-3,-2) .. controls (-2.333,0) and (-1.667,2) .. (-1,2)
            .. controls (-0.333,2) and (0.333,1.333) .. (1,0)
            .. controls (1.667,-1.333) and (2.333,-2) .. (3,-2);
  \pt{-3,-2}{left}{A}
  \pt{-1,2}{above}{B}
  \pt{1,0}{left}{C}
  \pt{3,-2}{below}{D}
  \node[coulCourbe,font=\footnotesize] at (-0.3,2.6) {$\Cf$};
\end{tikzpicture}
\end{center}
\explique{On place les quatre points, puis on trace un morceau de chaque tangente :}{la
pente donne le sens}
\explique{en $A(-3\,;-2)$, on avance de $1$ et on monte de $3$ ;}{$f'(-3)=3$}
\explique{en $B(-1\,;2)$ et en $D(3\,;-2)$, la tangente est horizontale ;}{$f'(-1)=0$ et
$f'(3)=0$}
\explique{en $C(1\,;0)$, on avance de $1$ et on descend de $2$.}{$f'(1)=-2$}
\explique{On trace enfin une courbe qui passe par les quatre points en \og
collant\fg{} à chaque tangente : elle monte jusqu'à $B$, puis descend jusqu'à
$D$.}{plusieurs courbes conviennent ; seuls les points et les tangentes sont
imposés}
\rem{Les tangentes d'abord, la courbe ensuite : là où la tangente est horizontale ($B$ et
$D$), la courbe forme un sommet en $B$ (maximum) ; en $D$, elle arrive à plat.}
\end{exo}

\partie{Équation de la tangente}

\begin{exo}[Avec un nombre dérivé admis]
\cas{a)}
\begin{calculs}
f(4) &= 4^{2}-3\times 4+2 &&\justif{on remplace $x$ par $4$}\\
f(4) &= 16-12+2 &&\\
f(4) &= 6 &&
\end{calculs}
\begin{calculs}
y &= f'(4)\,(x-4)+f(4) &&\justif{la formule, avec $a=4$}\\
y &= 5(x-4)+6 &&\justif{$f'(4)=5$ et $f(4)=6$}\\
y &= 5x-20+6 &&\justif{on développe}\\
y &= 5x-14 &&\justif{on réduit}
\end{calculs}
\explique{La tangente au point d'abscisse $4$ a pour équation $\res{y=5x-14}$.}{contrôle :
pour $x=4$, $y=20-14$, soit $6$ : c'est bien $f(4)$}
\cas{b)}
\begin{calculs}
f(1) &= 1^{2}-3\times 1+2 &&\justif{on remplace $x$ par $1$}\\
f(1) &= 0 &&
\end{calculs}
\begin{calculs}
y &= f'(1)\,(x-1)+f(1) &&\justif{la formule, avec $a=1$}\\
y &= -1(x-1)+0 &&\justif{$f'(1)=-1$ et $f(1)=0$}\\
y &= -x+1 &&\justif{on développe}
\end{calculs}
\explique{La tangente au point d'abscisse $1$ a pour équation $\res{y=-x+1}$.}{}
\end{exo}

\begin{exo}[Quatre tangentes]
\explique{Dans chaque cas, on recopie la formule $y=f'(a)\,(x-a)+f(a)$, on remplace,
puis on développe.}{propriété 1}
\cas{a)}
\begin{calculs}
y &= f'(2)\,(x-2)+f(2) &&\justif{la formule, avec $a=2$}\\
y &= 1(x-2)+3 &&\justif{$f'(2)=1$ et $f(2)=3$}\\
y &= x-2+3 &&\justif{on développe}\\
y &= \res{x+1} &&\justif{on réduit}
\end{calculs}
\cas{b)}
\begin{calculs}
y &= f'(-1)\,(x-(-1))+f(-1) &&\justif{la formule, avec $a=-1$}\\
y &= -2(x+1)+4 &&\justif{$x-(-1)=x+1$}\\
y &= -2x-2+4 &&\justif{on développe}\\
y &= \res{-2x+2} &&
\end{calculs}
\cas{c)}
\begin{calculs}
y &= f'(3)\,(x-3)+f(3) &&\justif{la formule, avec $a=3$}\\
y &= 0\times(x-3)+(-1) &&\justif{$f'(3)=0$ et $f(3)=-1$}\\
y &= \res{-1} &&\justif{tangente horizontale}
\end{calculs}
\cas{d)}
\begin{calculs}
y &= f'(0)\,(x-0)+f(0) &&\justif{la formule, avec $a=0$}\\
y &= \res{-3x+5} &&\justif{$x-0=x$}
\end{calculs}
\rem{Ne pas inverser les rôles : $f'(a)$, la pente, multiplie $(x-a)$ ; $f(a)$,
l'ordonnée du point, s'ajoute. Contrôle : pour $x=a$, on doit retrouver $y=f(a)$ ; en b),
$-2\times(-1)+2=4$.}
\end{exo}

\begin{exo}[Du taux de variation à la tangente]
\cas{a)}
\begin{calculs}
f(2) &= -2^{2}+3\times 2+1 &&\justif{$-2^{2}=-4$}\\
f(2) &= \res{3} &&\\[4pt]
f(2+h) &= -(2+h)^{2}+3(2+h)+1 &&\justif{on remplace $x$ par $(2+h)$}\\
f(2+h) &= -(4+4h+h^{2})+6+3h+1 &&\justif{identité remarquable}\\
f(2+h) &= -4-4h-h^{2}+6+3h+1 &&\justif{le $-$ change les trois signes}\\
f(2+h) &= \res{-h^{2}-h+3} &&\justif{on réduit}
\end{calculs}
\cas{b)}
\begin{calculs}
T(h) &= \frac{f(2+h)-f(2)}{h} &&\justif{la formule, avec $a=2$}\\
T(h) &= \frac{-h^{2}-h+3-3}{h} &&\justif{on remplace}\\
T(h) &= \frac{\barre{h}\,(-h-1)}{\barre{h}} &&\justif{on factorise par $h$, puis on simplifie}\\
T(h) &= -h-1 &&
\end{calculs}
\explique{Quand $h$ se rapproche de $0$, $T(h)=-h-1$ se rapproche de $-1$. Donc
$\res{f'(2)=-1}$.}{on remplace $h$ par $0$}
\cas{c)}
\begin{calculs}
y &= f'(2)\,(x-2)+f(2) &&\justif{la formule, avec $a=2$}\\
y &= -1(x-2)+3 &&\justif{$f'(2)=-1$ et $f(2)=3$}\\
y &= -x+2+3 &&\justif{on développe}\\
y &= -x+5 &&
\end{calculs}
\explique{La tangente au point d'abscisse $2$ a pour équation $\res{y=-x+5}$.}{}
\end{exo}

\partie{Approximation affine}

\begin{exo}[Le carré de $3{,}02$]
\cas{a)}
\begin{calculs}
f(3{,}02) &\approx 6\times 3{,}02-9 &&\justif{$3{,}02$ est proche de $3$ : on remplace $x$ dans l'équation de la tangente}\\
f(3{,}02) &\approx 18{,}12-9 &&\\
f(3{,}02) &\approx \res{9{,}12} &&
\end{calculs}
\begin{calculs}
f(3{,}02) &= 3{,}02^{2} &&\justif{on remplace $x$ par $3{,}02$ dans $f(x)$}\\
f(3{,}02) &= 9{,}1204 &&
\end{calculs}
\explique{L'écart entre $9{,}12$ et $9{,}1204$ n'est que de $0{,}0004$ :
l'approximation est très bonne.}{$3{,}02$ est très proche de $3$}
\cas{b)}
\begin{calculs}
f(2{,}99) &\approx 6\times 2{,}99-9 &&\justif{$2{,}99$ est proche de $3$}\\
f(2{,}99) &\approx 17{,}94-9 &&\\
f(2{,}99) &\approx \res{8{,}94} &&\\[4pt]
f(2{,}99) &= 2{,}99^{2} &&\justif{valeur exacte}\\
f(2{,}99) &= 8{,}9401 &&
\end{calculs}
\explique{L'écart n'est que de $0{,}0001$ : l'approximation est encore meilleure.}{$2{,}99$
est encore plus proche de $3$ que $3{,}02$}
\end{exo}

\begin{exo}[Une fonction inverse]
\cas{a)}
\begin{calculs}
f(2) &= \frac{1}{2-1} &&\justif{on remplace $x$ par $2$}\\
f(2) &= 1 &&\\[4pt]
y &= f'(2)\,(x-2)+f(2) &&\justif{la formule, avec $a=2$}\\
y &= -1(x-2)+1 &&\justif{$f'(2)=-1$ et $f(2)=1$}\\
y &= -x+2+1 &&\justif{on développe}\\
y &= -x+3 &&
\end{calculs}
\explique{La tangente au point d'abscisse $2$ a pour équation $\res{y=-x+3}$.}{}
\cas{b)}
\begin{calculs}
f(2{,}1) &\approx -2{,}1+3 &&\justif{$2{,}1$ est proche de $2$ : on remplace $x$ dans l'équation de la tangente}\\
f(2{,}1) &\approx \res{0{,}9} &&
\end{calculs}
\cas{c)}
\begin{calculs}
f(2{,}1) &= \frac{1}{2{,}1-1} &&\justif{on remplace $x$ par $2{,}1$ dans $f(x)$}\\
f(2{,}1) &= \frac{1}{1{,}1} &&\\
f(2{,}1) &\approx 0{,}909 &&\justif{à la calculatrice}
\end{calculs}
\explique{L'écart entre $0{,}9$ et $0{,}909$ est d'environ $0{,}009$ : l'approximation
est bonne.}{}
\end{exo}

\begin{exo}[Une plaque qui se dilate]
\cas{a)}
\begin{calculs}
A(10) &= 10^{2} &&\justif{aire de la plaque de $10$~cm de côté}\\
A(10) &= 100 &&\\[4pt]
y &= A'(10)\,(x-10)+A(10) &&\justif{la formule, avec $a=10$ ; ici, $x$ est le côté}\\
y &= 20(x-10)+100 &&\justif{$A'(10)=20$ et $A(10)=100$}\\
y &= 20x-200+100 &&\justif{on développe}\\
y &= 20x-100 &&
\end{calculs}
\explique{La tangente au point d'abscisse $10$ a pour équation $\res{y=20x-100}$.}{}
\cas{b)}
\begin{calculs}
A(10{,}1) &\approx 20\times 10{,}1-100 &&\justif{$10{,}1$ est proche de $10$}\\
A(10{,}1) &\approx 202-100 &&\\
A(10{,}1) &\approx 102 &&
\end{calculs}
\explique{L'aire de la plaque chauffée est d'environ $\res{102\text{ cm}^{2}}$.}{}
\cas{c)}
\begin{calculs}
A(10{,}1) &= 10{,}1^{2} &&\justif{valeur exacte}\\
A(10{,}1) &= 102{,}01 &&
\end{calculs}
\explique{L'aire exacte est $102{,}01$~cm$^{2}$ : l'approximation ne se trompe que de
$0{,}01$~cm$^{2}$.}{l'écart, $0{,}1^{2}$, est le petit carré du coin}
\end{exo}

\partie{Fonction dérivée d'une fonction polynôme}

\begin{exo}[Dériver]
\explique{On dérive terme par terme avec le tableau de la propriété 3 ; un nombre qui
multiplie reste devant.}{méthode 6}
\cas{a)}
\begin{calculs}
f'(x) &= \res{7} &&\justif{$mx+p$ devient $m$ : le $-5$ disparaît}
\end{calculs}
\cas{b)}
\begin{calculs}
g'(x) &= \res{-4} &&\justif{$1$ devient $0$ ; $-4x$ devient $-4$}
\end{calculs}
\cas{c)}
\begin{calculs}
h'(x) &= 2x+6+0 &&\justif{$x^{2}$ devient $2x$ ; $6x$ devient $6$ ; $-2$ devient $0$}\\
h'(x) &= \res{2x+6} &&
\end{calculs}
\cas{d)}
\begin{calculs}
k'(x) &= 3\times 2x-1+0 &&\justif{le $3$ reste devant ; $-x$ devient $-1$}\\
k'(x) &= \res{6x-1} &&
\end{calculs}
\cas{e)}
\begin{calculs}
m'(x) &= -5\times 2x+10 &&\justif{terme par terme}\\
m'(x) &= \res{-10x+10} &&
\end{calculs}
\cas{f)}
\begin{calculs}
p'(x) &= 3x^{2}+4\times 2x-1 &&\justif{$x^{3}$ devient $3x^{2}$}\\
p'(x) &= \res{3x^{2}+8x-1} &&
\end{calculs}
\cas{g)}
\begin{calculs}
q'(x) &= 2\times 4x^{3}-3\times 2x+0 &&\justif{$x^{4}$ devient $4x^{3}$}\\
q'(x) &= \res{8x^{3}-6x} &&
\end{calculs}
\cas{h)}
\begin{calculs}
r'(x) &= -5x^{4}+2\times 3x^{2} &&\justif{$x^{5}$ devient $5x^{4}$}\\
r'(x) &= \res{-5x^{4}+6x^{2}} &&
\end{calculs}
\rem{Un nombre seul (le $-5$ de $7x-5$, le $8$ de $3x^{2}-x+8$) a pour dérivée $0$ : il
disparaît.}
\end{exo}

\begin{exo}[Degré 3 et plus]
\explique{Pour $x^{n}$, l'exposant passe devant et on lui retire $1$.}{propriété 3}
\cas{a)}
\begin{calculs}
f'(x) &= 4\times 3x^{2}-6\times 2x+5-0 &&\justif{terme par terme}\\
f'(x) &= \res{12x^{2}-12x+5} &&
\end{calculs}
\cas{b)}
\begin{calculs}
g'(x) &= \res{6x^{5}-4x^{3}} &&\justif{$x^{6}$ devient $6x^{5}$ ; $x^{4}$ devient $4x^{3}$}
\end{calculs}
\cas{c)}
\begin{calculs}
h'(x) &= -2\times 5x^{4}+3x^{2}-7 &&\justif{$-7x$ devient $-7$}\\
h'(x) &= \res{-10x^{4}+3x^{2}-7} &&
\end{calculs}
\cas{d)}
\begin{calculs}
k'(x) &= 0-3\times 4x^{3} &&\justif{$5$ devient $0$}\\
k'(x) &= \res{-12x^{3}} &&
\end{calculs}
\cas{e)}
\begin{calculs}
m'(x) &= \frac{1}{3}\times 3x^{2}-2 &&\justif{le $\frac{1}{3}$ reste devant}\\
m'(x) &= \res{x^{2}-2} &&\justif{$\frac{1}{3}\times 3=1$}
\end{calculs}
\rem{Le nombre devant se garde : la dérivée de $-2x^{5}$ est $-2\times 5x^{4}=-10x^{4}$,
ni $5x^{4}$ ni $-2x^{4}$.}
\end{exo}

\begin{exo}[Entraînement]
\cas{a)}
\begin{calculs}
f'(x) &= \res{6} &&\justif{$mx+p$ devient $m$}
\end{calculs}
\cas{b)}
\begin{calculs}
g'(x) &= -2x+4-0 &&\justif{$-x^{2}$ devient $-2x$}\\
g'(x) &= \res{-2x+4} &&
\end{calculs}
\cas{c)}
\begin{calculs}
h'(x) &= 5\times 3x^{2}-2\times 2x &&\justif{terme par terme}\\
h'(x) &= \res{15x^{2}-4x} &&
\end{calculs}
\cas{d)}
\begin{calculs}
k'(x) &= \res{4x^{3}+3x^{2}+2x+1} &&\justif{chaque puissance perd $1$ ; $x$ devient $1$}
\end{calculs}
\cas{e)}
\begin{calculs}
m'(x) &= -3\times 4x^{3}+6 &&\\
m'(x) &= \res{-12x^{3}+6} &&
\end{calculs}
\cas{f)}
\begin{calculs}
p'(x) &= 2\times 5x^{4}-5\times 2x+0 &&\justif{$3$ devient $0$}\\
p'(x) &= \res{10x^{4}-10x} &&
\end{calculs}
\cas{g)}
\begin{calculs}
q'(x) &= \frac{1}{2}\times 2x+3 &&\justif{le $\frac{1}{2}$ reste devant}\\
q'(x) &= \res{x+3} &&\justif{$\frac{1}{2}\times 2=1$}
\end{calculs}
\cas{h)}
\begin{calculs}
r'(x) &= \frac{1}{4}\times 4x^{3}-1 &&\\
r'(x) &= \res{x^{3}-1} &&\justif{$\frac{1}{4}\times 4=1$}
\end{calculs}
\end{exo}

\begin{exo}[Développer d'abord]
\explique{Le tableau de la propriété 3 ne donne la dérivée ni d'un produit ni d'un
carré : on développe d'abord, pour obtenir un polynôme.}{}
\cas{a)}
\begin{calculs}
A(x) &= (x-1)(x+4) &&\justif{on recopie l'expression}\\
A(x) &= x^{2}+4x-x-4 &&\justif{double distributivité}\\
A(x) &= x^{2}+3x-4 &&\justif{on réduit}\\
A'(x) &= \res{2x+3} &&\justif{terme par terme}
\end{calculs}
\cas{b)}
\begin{calculs}
B(x) &= (2x+1)(x-3) &&\\
B(x) &= 2x^{2}-6x+x-3 &&\justif{double distributivité}\\
B(x) &= 2x^{2}-5x-3 &&\\
B'(x) &= \res{4x-5} &&\justif{$2\times 2x=4x$}
\end{calculs}
\cas{c)}
\begin{calculs}
C(x) &= (x+2)^{2} &&\\
C(x) &= x^{2}+4x+4 &&\justif{$(a+b)^{2}=a^{2}+2ab+b^{2}$}\\
C'(x) &= \res{2x+4} &&
\end{calculs}
\cas{d)}
\begin{calculs}
D(x) &= x(x^{2}-5) &&\\
D(x) &= x^{3}-5x &&\justif{on distribue $x$}\\
D'(x) &= \res{3x^{2}-5} &&
\end{calculs}
\rem{Erreur fréquente : dériver chaque facteur et multiplier, $(x-1)'\times(x+4)'=1$ au
lieu de $2x+3$.}
\end{exo}

\begin{exo}[Dérivée et tangentes]
\cas{a)}
\begin{calculs}
f'(x) &= 3x^{2}-2+0 &&\justif{$x^{3}$ devient $3x^{2}$ ; $-2x$ devient $-2$}\\
f'(x) &= \res{3x^{2}-2} &&
\end{calculs}
\cas{b)}
\begin{calculs}
f(1) &= 1^{3}-2\times 1+1 &&\justif{on remplace $x$ par $1$ dans $f(x)$}\\
f(1) &= 0 &&\\
f'(1) &= 3\times 1^{2}-2 &&\justif{on remplace $x$ par $1$ dans $f'(x)$}\\
f'(1) &= 1 &&
\end{calculs}
\begin{calculs}
y &= f'(1)\,(x-1)+f(1) &&\justif{la formule, avec $a=1$}\\
y &= 1(x-1)+0 &&\\
y &= x-1 &&
\end{calculs}
\explique{La tangente au point d'abscisse $1$ a pour équation $\res{y=x-1}$.}{}
\cas{c)}
\begin{calculs}
f(0) &= 0-0+1 &&\\
f(0) &= 1 &&\\
f'(0) &= 3\times 0^{2}-2 &&\\
f'(0) &= -2 &&\\[4pt]
y &= f'(0)\,(x-0)+f(0) &&\justif{la formule, avec $a=0$}\\
y &= -2x+1 &&
\end{calculs}
\explique{La tangente au point d'abscisse $0$ a pour équation $\res{y=-2x+1}$.}{}
\rem{Ne pas confondre $f(1)$, calculé avec $f(x)$, et $f'(1)$, calculé avec $f'(x)$.}
\end{exo}

\partie{Autres fonctions dérivées}

\begin{exo}[Inverse, cosinus, sinus]
\explique{$\frac{1}{x}$ devient $-\frac{1}{x^{2}}$, $\cos x$ devient $-\sin x$, $\sin x$
devient $\cos x$ ; le nombre qui multiplie reste devant.}{propriété 4}
\cas{a)}
\begin{calculs}
f'(x) &= 4\times\left(-\frac{1}{x^{2}}\right) &&\justif{$\frac{4}{x}=4\times\frac{1}{x}$ : le $4$ reste devant}\\
f'(x) &= \res{-\frac{4}{x^{2}}} &&
\end{calculs}
\cas{b)}
\begin{calculs}
g'(x) &= -7\times\left(-\frac{1}{x^{2}}\right) &&\justif{le $-7$ reste devant}\\
g'(x) &= \res{\frac{7}{x^{2}}} &&\justif{moins par moins : plus}
\end{calculs}
\cas{c)}
\begin{calculs}
h'(x) &= \res{-\sin x+3} &&\justif{$\cos x$ devient $-\sin x$ ; $3x$ devient $3$}
\end{calculs}
\cas{d)}
\begin{calculs}
k'(x) &= \res{5\cos x} &&\justif{$\sin x$ devient $\cos x$, le $5$ reste devant}
\end{calculs}
\cas{e)}
\begin{calculs}
m'(x) &= 2\times(-\sin x)-3\cos x &&\justif{terme par terme}\\
m'(x) &= \res{-2\sin x-3\cos x} &&
\end{calculs}
\cas{f)}
\begin{calculs}
p'(x) &= 4x^{3}-\left(-\frac{1}{x^{2}}\right) &&\justif{$x^{4}$ devient $4x^{3}$}\\
p'(x) &= \res{4x^{3}+\frac{1}{x^{2}}} &&
\end{calculs}
\cas{g)}
\begin{calculs}
q'(x) &= 6\times\left(-\frac{1}{x^{2}}\right)+2\times 2x &&\\
q'(x) &= \res{-\frac{6}{x^{2}}+4x} &&
\end{calculs}
\cas{h)}
\begin{calculs}
r'(x) &= \res{-\cos x+3x^{2}} &&\justif{$-\sin x$ devient $-\cos x$}
\end{calculs}
\rem{Seule la dérivée de $\cos x$ porte un signe $-$ : $(\cos x)'=-\sin x$, mais
$(\sin x)'=\cos x$.}
\end{exo}

\begin{exo}[Entraînement]
\cas{a)}
\begin{calculs}
f'(x) &= 3\times\left(-\frac{1}{x^{2}}\right)+1 &&\justif{$x$ devient $1$}\\
f'(x) &= \res{-\frac{3}{x^{2}}+1} &&
\end{calculs}
\cas{b)}
\begin{calculs}
g'(x) &= 4\times(-\sin x)+2x &&\\
g'(x) &= \res{-4\sin x+2x} &&
\end{calculs}
\cas{c)}
\begin{calculs}
h'(x) &= \res{-2\cos x} &&\justif{$\sin x$ devient $\cos x$ ; $5$ devient $0$}
\end{calculs}
\cas{d)}
\begin{calculs}
k'(x) &= 3x^{2}-5\times\left(-\frac{1}{x^{2}}\right) &&\\
k'(x) &= \res{3x^{2}+\frac{5}{x^{2}}} &&\justif{moins par moins : plus}
\end{calculs}
\cas{e)}
\begin{calculs}
m'(x) &= \cos x-(-\sin x) &&\justif{$-\cos x$ devient $-(-\sin x)$}\\
m'(x) &= \res{\cos x+\sin x} &&
\end{calculs}
\cas{f)}
\begin{calculs}
p'(x) &= -\frac{1}{x^{2}}-3\times(-\sin x) &&\\
p'(x) &= \res{-\frac{1}{x^{2}}+3\sin x} &&
\end{calculs}
\end{exo}

\begin{exo}[Tangentes à la courbe de $\frac{2}{x}$]
\cas{a)}
\begin{calculs}
f'(x) &= 2\times\left(-\frac{1}{x^{2}}\right) &&\justif{le $2$ reste devant}\\
f'(x) &= \res{-\frac{2}{x^{2}}} &&
\end{calculs}
\cas{b)}
\begin{calculs}
f(1) &= \frac{2}{1} = 2 &&\justif{on remplace $x$ par $1$ dans $f(x)$}\\
f'(1) &= -\frac{2}{1^{2}} = -2 &&\justif{puis dans $f'(x)$}\\[4pt]
y &= f'(1)\,(x-1)+f(1) &&\justif{la formule, avec $a=1$}\\
y &= -2(x-1)+2 &&\\
y &= -2x+2+2 &&\justif{on développe}\\
y &= -2x+4 &&
\end{calculs}
\explique{La tangente au point d'abscisse $1$ a pour équation $\res{y=-2x+4}$.}{}
\cas{c)}
\begin{calculs}
f(2) &= \frac{2}{2} = 1 &&\\
f'(2) &= -\frac{2}{2^{2}} = -\frac{1}{2} &&\justif{$\frac{2}{4}=\frac{1}{2}$}\\[4pt]
y &= f'(2)\,(x-2)+f(2) &&\justif{la formule, avec $a=2$}\\
y &= -\frac{1}{2}(x-2)+1 &&\\
y &= -\frac{1}{2}x+1+1 &&\justif{$-\frac{1}{2}\times(-2)=1$}\\
y &= -\frac{1}{2}x+2 &&
\end{calculs}
\explique{La tangente au point d'abscisse $2$ a pour équation
$\res{y=-\frac{1}{2}x+2}$.}{}
\end{exo}

\begin{exo}[Sinus et cosinus près de $0$]
\cas{a)}
\begin{calculs}
f'(x) &= \cos x &&\justif{$\sin x$ devient $\cos x$}\\
f'(0) &= \cos 0 = 1 &&\\
f(0) &= \sin 0 = 0 &&\\[4pt]
y &= f'(0)\,(x-0)+f(0) &&\justif{la formule, avec $a=0$}\\
y &= 1\times x+0 &&\\
y &= x &&
\end{calculs}
\explique{La tangente au point d'abscisse $0$ a pour équation $\res{y=x}$.}{}
\cas{b)}
\begin{calculs}
\sin(0{,}1) &\approx 0{,}1 &&\justif{$0{,}1$ est proche de $0$ : on remplace $x$ dans $y=x$}
\end{calculs}
\explique{La calculatrice, en radians, donne $\sin(0{,}1)\approx 0{,}0998$ : l'écart est
d'environ $0{,}0002$. Donc $\sin(0{,}1)\approx\res{0{,}1}$.}{en degrés, on obtiendrait
une tout autre valeur}
\cas{c)}
\begin{calculs}
g'(x) &= -\sin x &&\justif{$\cos x$ devient $-\sin x$}\\
g'(0) &= -\sin 0 &&\\
g'(0) &= \res{0} &&
\end{calculs}
\explique{Le nombre dérivé est nul : la tangente au point d'abscisse $0$ est
\res{horizontale}. Comme $g(0)=\cos 0=1$, elle a pour équation $y=1$.}{près de $0$,
$\cos x\approx 1$}
\end{exo}

\partie{Dérivée d'un produit, d'un carré}

\begin{exo}[Des carrés]
\explique{On écrit $u$ et $u'$, puis on recopie la formule $2\times u'\times u$ et on
remplace.}{méthode 8}
\cas{a)}
\begin{calculs}
u(x) &= 4x+1 &&\justif{donc $u'(x)=4$}\\[4pt]
f'(x) &= 2\times u'(x)\times u(x) &&\justif{la formule du carré}\\
f'(x) &= 2\times 4\times(4x+1) &&\justif{on remplace}\\
f'(x) &= \res{32x+8} &&\justif{on développe $8(4x+1)$}
\end{calculs}
\cas{b)}
\begin{calculs}
u(x) &= 2-5x &&\justif{donc $u'(x)=-5$}\\[4pt]
g'(x) &= 2\times u'(x)\times u(x) &&\\
g'(x) &= 2\times(-5)\times(2-5x) &&\justif{on remplace}\\
g'(x) &= -10(2-5x) &&\\
g'(x) &= \res{50x-20} &&\justif{$-10\times(-5x)=+50x$}
\end{calculs}
\cas{c)}
\begin{calculs}
u(x) &= x^{2}+3 &&\justif{donc $u'(x)=2x$}\\[4pt]
h'(x) &= 2\times u'(x)\times u(x) &&\\
h'(x) &= 2\times 2x\times(x^{2}+3) &&\\
h'(x) &= 4x(x^{2}+3) &&\\
h'(x) &= \res{4x^{3}+12x} &&
\end{calculs}
\cas{d)}
\begin{calculs}
u(x) &= x^{3}-1 &&\justif{donc $u'(x)=3x^{2}$}\\[4pt]
k'(x) &= 2\times u'(x)\times u(x) &&\\
k'(x) &= 2\times 3x^{2}\times(x^{3}-1) &&\\
k'(x) &= 6x^{2}(x^{3}-1) &&\\
k'(x) &= \res{6x^{5}-6x^{2}} &&\justif{$x^{2}\times x^{3}=x^{5}$}
\end{calculs}
\rem{Erreur fréquente : $\big((4x+1)^{2}\big)'=2(4x+1)$, en oubliant de multiplier par
$u'=4$.}
\end{exo}

\begin{exo}[Des produits]
\explique{On nomme les deux facteurs $u$ et $v$, on calcule $u'$ et $v'$, puis on recopie
la formule $u'\times v+u\times v'$ et on remplace.}{méthode 8}
\cas{a)}
\begin{calculs}
u(x) &= x+2 &&\justif{donc $u'(x)=1$}\\
v(x) &= 3x-1 &&\justif{donc $v'(x)=3$}\\[4pt]
f'(x) &= u'(x)\,v(x)+u(x)\,v'(x) &&\justif{la formule du produit}\\
f'(x) &= 1\times(3x-1)+(x+2)\times 3 &&\justif{on remplace}\\
f'(x) &= 3x-1+3x+6 &&\justif{on développe}\\
f'(x) &= \res{6x+5} &&\justif{on réduit}
\end{calculs}
\cas{b)}
\begin{calculs}
u(x) &= 2x-3 &&\justif{donc $u'(x)=2$}\\
v(x) &= x^{2}+1 &&\justif{donc $v'(x)=2x$}\\[4pt]
g'(x) &= u'(x)\,v(x)+u(x)\,v'(x) &&\\
g'(x) &= 2(x^{2}+1)+(2x-3)\times 2x &&\\
g'(x) &= 2x^{2}+2+4x^{2}-6x &&\\
g'(x) &= \res{6x^{2}-6x+2} &&
\end{calculs}
\cas{c)}
\begin{calculs}
u(x) &= x^{2} &&\justif{donc $u'(x)=2x$}\\
v(x) &= 4x+5 &&\justif{donc $v'(x)=4$}\\[4pt]
h'(x) &= u'(x)\,v(x)+u(x)\,v'(x) &&\\
h'(x) &= 2x(4x+5)+x^{2}\times 4 &&\\
h'(x) &= 8x^{2}+10x+4x^{2} &&\\
h'(x) &= \res{12x^{2}+10x} &&
\end{calculs}
\cas{d)}
\begin{calculs}
u(x) &= x^{2}-x &&\justif{donc $u'(x)=2x-1$}\\
v(x) &= x^{2}+2 &&\justif{donc $v'(x)=2x$}\\[4pt]
k'(x) &= u'(x)\,v(x)+u(x)\,v'(x) &&\\
k'(x) &= (2x-1)(x^{2}+2)+(x^{2}-x)\times 2x &&\\
k'(x) &= 2x^{3}+4x-x^{2}-2+2x^{3}-2x^{2} &&\justif{double distributivité}\\
k'(x) &= \res{4x^{3}-3x^{2}+4x-2} &&
\end{calculs}
\cas{e)}
\begin{calculs}
u(x) &= x &&\justif{donc $u'(x)=1$}\\
v(x) &= \sin x &&\justif{donc $v'(x)=\cos x$}\\[4pt]
m'(x) &= u'(x)\,v(x)+u(x)\,v'(x) &&\\
m'(x) &= 1\times\sin x+x\times\cos x &&\\
m'(x) &= \res{\sin x+x\cos x} &&
\end{calculs}
\rem{La dérivée d'un produit n'est pas le produit des dérivées : en a), $1\times 3=3$
serait faux.}
\end{exo}

\begin{exo}[Entraînement]
\cas{a)}
\begin{calculs}
u(x) &= 3x+1 &&\justif{donc $u'(x)=3$}\\
v(x) &= x-2 &&\justif{donc $v'(x)=1$}\\[4pt]
f'(x) &= u'(x)\,v(x)+u(x)\,v'(x) &&\justif{la formule du produit}\\
f'(x) &= 3(x-2)+(3x+1)\times 1 &&\\
f'(x) &= 3x-6+3x+1 &&\\
f'(x) &= \res{6x-5} &&
\end{calculs}
\cas{b)}
\begin{calculs}
u(x) &= x+5 &&\justif{donc $u'(x)=1$}\\[4pt]
g'(x) &= 2\times u'(x)\times u(x) &&\justif{la formule du carré}\\
g'(x) &= 2\times 1\times(x+5) &&\\
g'(x) &= \res{2x+10} &&
\end{calculs}
\cas{c)}
\begin{calculs}
u(x) &= x^{2}-4 &&\justif{donc $u'(x)=2x$}\\
v(x) &= 2x+3 &&\justif{donc $v'(x)=2$}\\[4pt]
h'(x) &= u'(x)\,v(x)+u(x)\,v'(x) &&\\
h'(x) &= 2x(2x+3)+(x^{2}-4)\times 2 &&\\
h'(x) &= 4x^{2}+6x+2x^{2}-8 &&\\
h'(x) &= \res{6x^{2}+6x-8} &&
\end{calculs}
\cas{d)}
\begin{calculs}
u(x) &= 3-x^{2} &&\justif{donc $u'(x)=-2x$}\\[4pt]
k'(x) &= 2\times u'(x)\times u(x) &&\justif{la formule du carré}\\
k'(x) &= 2\times(-2x)\times(3-x^{2}) &&\\
k'(x) &= -4x(3-x^{2}) &&\\
k'(x) &= \res{4x^{3}-12x} &&\justif{$-4x\times(-x^{2})=+4x^{3}$}
\end{calculs}
\cas{e)}
\begin{calculs}
u(x) &= x^{2} &&\justif{donc $u'(x)=2x$}\\
v(x) &= \cos x &&\justif{donc $v'(x)=-\sin x$}\\[4pt]
m'(x) &= u'(x)\,v(x)+u(x)\,v'(x) &&\\
m'(x) &= 2x\cos x+x^{2}\times(-\sin x) &&\\
m'(x) &= \res{2x\cos x-x^{2}\sin x} &&
\end{calculs}
\cas{f)}
\begin{calculs}
u(x) &= 2x-1 &&\justif{donc $u'(x)=2$}\\
v(x) &= x^{3}+x &&\justif{donc $v'(x)=3x^{2}+1$}\\[4pt]
p'(x) &= u'(x)\,v(x)+u(x)\,v'(x) &&\\
p'(x) &= 2(x^{3}+x)+(2x-1)(3x^{2}+1) &&\\
p'(x) &= 2x^{3}+2x+6x^{3}+2x-3x^{2}-1 &&\justif{double distributivité}\\
p'(x) &= \res{8x^{3}-3x^{2}+4x-1} &&
\end{calculs}
\end{exo}

\begin{exo}[Deux façons de faire]
\cas{a)}
\begin{calculs}
u(x) &= x+3 &&\justif{donc $u'(x)=1$}\\
v(x) &= 2x-1 &&\justif{donc $v'(x)=2$}\\[4pt]
f'(x) &= u'(x)\,v(x)+u(x)\,v'(x) &&\justif{la formule du produit}\\
f'(x) &= 1\times(2x-1)+(x+3)\times 2 &&\\
f'(x) &= 2x-1+2x+6 &&\\
f'(x) &= \res{4x+5} &&
\end{calculs}
\cas{b)}
\begin{calculs}
f(x) &= (x+3)(2x-1) &&\justif{on recopie l'expression}\\
f(x) &= 2x^{2}-x+6x-3 &&\justif{double distributivité}\\
f(x) &= 2x^{2}+5x-3 &&\\
f'(x) &= 2\times 2x+5 &&\justif{terme par terme}\\
f'(x) &= \res{4x+5} &&
\end{calculs}
\cas{c)}
\explique{Les deux méthodes donnent le même résultat, $4x+5$.}{développer permet de
vérifier une dérivée de produit}
\rem{La dérivée d'un produit n'est pas le produit des dérivées : $1\times 2=2$ serait
faux.}
\end{exo}

\partie{Dérivée de l'inverse, d'un quotient}

\begin{exo}[Des inverses]
\explique{On écrit $v$ et $v'$, puis on recopie la formule $-\dfrac{v'}{v^{2}}$ et on
remplace.}{méthode 9}
\cas{a)}
\begin{calculs}
v(x) &= 3x+2 &&\justif{donc $v'(x)=3$}\\[4pt]
f'(x) &= -\frac{v'(x)}{v(x)^{2}} &&\justif{la formule de l'inverse}\\
f'(x) &= \res{-\frac{3}{(3x+2)^{2}}} &&\justif{on remplace}
\end{calculs}
\cas{b)}
\begin{calculs}
v(x) &= 1-4x &&\justif{donc $v'(x)=-4$}\\[4pt]
g'(x) &= -\frac{v'(x)}{v(x)^{2}} &&\\
g'(x) &= -\frac{-4}{(1-4x)^{2}} &&\justif{on remplace}\\
g'(x) &= \res{\frac{4}{(1-4x)^{2}}} &&\justif{moins par moins : plus}
\end{calculs}
\cas{c)}
\begin{calculs}
v(x) &= x^{2}+1 &&\justif{donc $v'(x)=2x$}\\[4pt]
h'(x) &= -\frac{v'(x)}{v(x)^{2}} &&\\
h'(x) &= \res{-\frac{2x}{(x^{2}+1)^{2}}} &&
\end{calculs}
\cas{d)}
\explique{$k(x)=5\times\dfrac{1}{2x-1}$ : le $5$ reste devant.}{comme $\frac{5}{x}=5\times\frac{1}{x}$}\nopagebreak
\begin{calculs}
v(x) &= 2x-1 &&\justif{donc $v'(x)=2$}\\[4pt]
k'(x) &= 5\times\left(-\frac{v'(x)}{v(x)^{2}}\right) &&\\
k'(x) &= 5\times\left(-\frac{2}{(2x-1)^{2}}\right) &&\\
k'(x) &= \res{-\frac{10}{(2x-1)^{2}}} &&
\end{calculs}
\rem{Le dénominateur $v^{2}$ reste écrit au carré : on ne le développe pas.}
\end{exo}

\begin{exo}[Des quotients]
\explique{On écrit $u$, $u'$, $v$ et $v'$, puis on recopie la formule et on
remplace.}{$\dfrac{u'\times v-u\times v'}{v^{2}}$}
\cas{a)}
\begin{calculs}
u(x) &= x+4 &&\justif{donc $u'(x)=1$}\\
v(x) &= x-1 &&\justif{donc $v'(x)=1$}\\[4pt]
f'(x) &= \frac{u'(x)\,v(x)-u(x)\,v'(x)}{v(x)^{2}} &&\justif{la formule du quotient}\\
f'(x) &= \frac{1\times(x-1)-(x+4)\times 1}{(x-1)^{2}} &&\justif{parenthèses autour de $u(x)$}\\
f'(x) &= \frac{x-1-x-4}{(x-1)^{2}} &&\justif{le $-$ porte sur $x$ et sur $4$}\\
f'(x) &= \res{\frac{-5}{(x-1)^{2}}} &&\justif{on réduit}
\end{calculs}
\cas{b)}
\begin{calculs}
u(x) &= 2x-1 &&\justif{donc $u'(x)=2$}\\
v(x) &= x+2 &&\justif{donc $v'(x)=1$}\\[4pt]
g'(x) &= \frac{u'(x)\,v(x)-u(x)\,v'(x)}{v(x)^{2}} &&\\
g'(x) &= \frac{2(x+2)-(2x-1)\times 1}{(x+2)^{2}} &&\\
g'(x) &= \frac{2x+4-2x+1}{(x+2)^{2}} &&\justif{$-(-1)=+1$}\\
g'(x) &= \res{\frac{5}{(x+2)^{2}}} &&
\end{calculs}
\cas{c)}
\begin{calculs}
u(x) &= 3x &&\justif{donc $u'(x)=3$}\\
v(x) &= x+1 &&\justif{donc $v'(x)=1$}\\[4pt]
h'(x) &= \frac{u'(x)\,v(x)-u(x)\,v'(x)}{v(x)^{2}} &&\\
h'(x) &= \frac{3(x+1)-3x\times 1}{(x+1)^{2}} &&\\
h'(x) &= \frac{3x+3-3x}{(x+1)^{2}} &&\\
h'(x) &= \res{\frac{3}{(x+1)^{2}}} &&
\end{calculs}
\cas{d)}
\begin{calculs}
u(x) &= x^{2}+1 &&\justif{donc $u'(x)=2x$}\\
v(x) &= x-2 &&\justif{donc $v'(x)=1$}\\[4pt]
k'(x) &= \frac{u'(x)\,v(x)-u(x)\,v'(x)}{v(x)^{2}} &&\\
k'(x) &= \frac{2x(x-2)-(x^{2}+1)\times 1}{(x-2)^{2}} &&\\
k'(x) &= \frac{2x^{2}-4x-x^{2}-1}{(x-2)^{2}} &&\justif{on développe le numérateur}\\
k'(x) &= \res{\frac{x^{2}-4x-1}{(x-2)^{2}}} &&
\end{calculs}
\rem{Des parenthèses autour de $u(x)$ : le signe $-$ porte sur tout le terme. En a),
$-(x+4)=-x-4$, et non $-x+4$.}
\end{exo}

\begin{exo}[Entraînement]
\cas{a)}
\begin{calculs}
v(x) &= x+5 &&\justif{donc $v'(x)=1$}\\[4pt]
f'(x) &= -\frac{v'(x)}{v(x)^{2}} &&\justif{la formule de l'inverse}\\
f'(x) &= \res{-\frac{1}{(x+5)^{2}}} &&
\end{calculs}
\cas{b)}
\explique{$g(x)=3\times\dfrac{1}{4x-1}$ : le $3$ reste devant.}{}\nopagebreak
\begin{calculs}
v(x) &= 4x-1 &&\justif{donc $v'(x)=4$}\\[4pt]
g'(x) &= 3\times\left(-\frac{4}{(4x-1)^{2}}\right) &&\\
g'(x) &= \res{-\frac{12}{(4x-1)^{2}}} &&
\end{calculs}
\cas{c)}
\begin{calculs}
u(x) &= x-3 &&\justif{donc $u'(x)=1$}\\
v(x) &= x+2 &&\justif{donc $v'(x)=1$}\\[4pt]
h'(x) &= \frac{u'(x)\,v(x)-u(x)\,v'(x)}{v(x)^{2}} &&\justif{la formule du quotient}\\
h'(x) &= \frac{1\times(x+2)-(x-3)\times 1}{(x+2)^{2}} &&\\
h'(x) &= \frac{x+2-x+3}{(x+2)^{2}} &&\justif{$-(-3)=+3$}\\
h'(x) &= \res{\frac{5}{(x+2)^{2}}} &&
\end{calculs}
\cas{d)}
\begin{calculs}
u(x) &= 4x+1 &&\justif{donc $u'(x)=4$}\\
v(x) &= 2x-3 &&\justif{donc $v'(x)=2$}\\[4pt]
k'(x) &= \frac{u'(x)\,v(x)-u(x)\,v'(x)}{v(x)^{2}} &&\\
k'(x) &= \frac{4(2x-3)-(4x+1)\times 2}{(2x-3)^{2}} &&\\
k'(x) &= \frac{8x-12-8x-2}{(2x-3)^{2}} &&\\
k'(x) &= \res{\frac{-14}{(2x-3)^{2}}} &&
\end{calculs}
\cas{e)}
\explique{$m(x)=2\times\dfrac{1}{x^{2}+4}$ : le $2$ reste devant.}{}\nopagebreak
\begin{calculs}
v(x) &= x^{2}+4 &&\justif{donc $v'(x)=2x$}\\[4pt]
m'(x) &= 2\times\left(-\frac{2x}{(x^{2}+4)^{2}}\right) &&\\
m'(x) &= \res{-\frac{4x}{(x^{2}+4)^{2}}} &&
\end{calculs}
\cas{f)}
\begin{calculs}
u(x) &= x^{2} &&\justif{donc $u'(x)=2x$}\\
v(x) &= x+1 &&\justif{donc $v'(x)=1$}\\[4pt]
p'(x) &= \frac{u'(x)\,v(x)-u(x)\,v'(x)}{v(x)^{2}} &&\\
p'(x) &= \frac{2x(x+1)-x^{2}\times 1}{(x+1)^{2}} &&\\
p'(x) &= \frac{2x^{2}+2x-x^{2}}{(x+1)^{2}} &&\\
p'(x) &= \res{\frac{x^{2}+2x}{(x+1)^{2}}} &&
\end{calculs}
\rem{Un nombre seul au numérateur (b et e) : la formule de l'inverse suffit, le nombre
reste devant.}
\end{exo}

\begin{exo}[Quotient et tangente]
\cas{a)}
\begin{calculs}
u(x) &= 2x+1 &&\justif{donc $u'(x)=2$}\\
v(x) &= x-1 &&\justif{donc $v'(x)=1$}\\[4pt]
f'(x) &= \frac{u'(x)\,v(x)-u(x)\,v'(x)}{v(x)^{2}} &&\justif{la formule du quotient}\\
f'(x) &= \frac{2(x-1)-(2x+1)\times 1}{(x-1)^{2}} &&\\
f'(x) &= \frac{2x-2-2x-1}{(x-1)^{2}} &&\\
f'(x) &= \res{\frac{-3}{(x-1)^{2}}} &&
\end{calculs}
\cas{b)}
\begin{calculs}
f(2) &= \frac{2\times 2+1}{2-1} = 5 &&\justif{on remplace $x$ par $2$ dans $f(x)$}\\
f'(2) &= \frac{-3}{(2-1)^{2}} = -3 &&\justif{puis dans $f'(x)$}\\[4pt]
y &= f'(2)\,(x-2)+f(2) &&\justif{la formule, avec $a=2$}\\
y &= -3(x-2)+5 &&\\
y &= -3x+6+5 &&\justif{on développe}\\
y &= -3x+11 &&
\end{calculs}
\explique{La tangente au point d'abscisse $2$ a pour équation $\res{y=-3x+11}$.}{}
\end{exo}

\partie{Fonctions composées}

\begin{exo}[Des puissances]
\explique{On écrit $u$, $u'$ et $n$, puis on recopie la formule
$n\times u'\times u^{n-1}$ et on remplace ; on ne développe pas.}{méthode 10}
\cas{a)}
\begin{calculs}
u(x) &= 3x-2 &&\justif{donc $u'(x)=3$ ; ici $n=4$}\\[4pt]
f'(x) &= 4\times u'(x)\times u(x)^{3} &&\justif{la formule de $u^{n}$}\\
f'(x) &= 4\times 3\times(3x-2)^{3} &&\justif{on remplace}\\
f'(x) &= \res{12(3x-2)^{3}} &&
\end{calculs}
\cas{b)}
\begin{calculs}
u(x) &= x^{2}+1 &&\justif{donc $u'(x)=2x$ ; ici $n=3$}\\[4pt]
g'(x) &= 3\times u'(x)\times u(x)^{2} &&\\
g'(x) &= 3\times 2x\times(x^{2}+1)^{2} &&\\
g'(x) &= \res{6x(x^{2}+1)^{2}} &&
\end{calculs}
\cas{c)}
\begin{calculs}
u(x) &= 1-2x &&\justif{donc $u'(x)=-2$ ; ici $n=5$}\\[4pt]
h'(x) &= 5\times u'(x)\times u(x)^{4} &&\\
h'(x) &= 5\times(-2)\times(1-2x)^{4} &&\\
h'(x) &= \res{-10(1-2x)^{4}} &&
\end{calculs}
\cas{d)}
\begin{calculs}
u(x) &= x^{3}+x &&\justif{donc $u'(x)=3x^{2}+1$ ; ici $n=2$}\\[4pt]
k'(x) &= 2\times u'(x)\times u(x)^{1} &&\justif{$u^{1}=u$ : c'est la formule du carré}\\
k'(x) &= \res{2(3x^{2}+1)(x^{3}+x)} &&
\end{calculs}
\rem{Erreur fréquente : oublier $u'$. En a), $4(3x-2)^{3}$ est faux : il manque le
facteur $3$.}
\end{exo}

\begin{exo}[Cosinus et sinus composés]
\explique{On écrit $u$ et $u'$ ; $\cos(u)$ devient $-u'\times\sin(u)$, $\sin(u)$ devient
$u'\times\cos(u)$.}{propriété 7}
\cas{a)}
\begin{calculs}
u(x) &= 4x &&\justif{donc $u'(x)=4$}\\[4pt]
f'(x) &= -u'(x)\times\sin\big(u(x)\big) &&\justif{la formule de $\cos(u)$}\\
f'(x) &= \res{-4\sin(4x)} &&\justif{on remplace}
\end{calculs}
\cas{b)}
\begin{calculs}
u(x) &= 2x+3 &&\justif{donc $u'(x)=2$}\\[4pt]
g'(x) &= u'(x)\times\cos\big(u(x)\big) &&\justif{la formule de $\sin(u)$}\\
g'(x) &= \res{2\cos(2x+3)} &&
\end{calculs}
\cas{c)}
\begin{calculs}
u(x) &= 5x-1 &&\justif{donc $u'(x)=5$}\\[4pt]
h'(x) &= 3\times\big(-5\sin(5x-1)\big) &&\justif{la formule de $\cos(u)$ ; le $3$ reste devant}\\
h'(x) &= \res{-15\sin(5x-1)} &&
\end{calculs}
\cas{d)}
\begin{calculs}
u(t) &= 100t &&\justif{la variable est $t$ ; $u'(t)=100$}\\[4pt]
k'(t) &= 2\times 100\cos(100t) &&\justif{la formule de $\sin(u)$ ; le $2$ reste devant}\\
k'(t) &= \res{200\cos(100t)} &&
\end{calculs}
\end{exo}

\begin{exo}[Entraînement]
\cas{a)}
\begin{calculs}
u(x) &= 5x+2 &&\justif{donc $u'(x)=5$ ; ici $n=3$}\\[4pt]
f'(x) &= 3\times u'(x)\times u(x)^{2} &&\justif{la formule de $u^{n}$}\\
f'(x) &= 3\times 5\times(5x+2)^{2} &&\\
f'(x) &= \res{15(5x+2)^{2}} &&
\end{calculs}
\cas{b)}
\begin{calculs}
u(x) &= x^{2}-3x &&\justif{donc $u'(x)=2x-3$ ; ici $n=4$}\\[4pt]
g'(x) &= 4\times u'(x)\times u(x)^{3} &&\\
g'(x) &= \res{4(2x-3)(x^{2}-3x)^{3}} &&
\end{calculs}
\cas{c)}
\begin{calculs}
u(x) &= 3x-2 &&\justif{donc $u'(x)=3$}\\[4pt]
h'(x) &= -u'(x)\times\sin\big(u(x)\big) &&\justif{la formule de $\cos(u)$}\\
h'(x) &= \res{-3\sin(3x-2)} &&
\end{calculs}
\cas{d)}
\begin{calculs}
u(x) &= x^{2} &&\justif{donc $u'(x)=2x$}\\[4pt]
k'(x) &= 4\times 2x\cos(x^{2}) &&\justif{la formule de $\sin(u)$ ; le $4$ reste devant}\\
k'(x) &= \res{8x\cos(x^{2})} &&
\end{calculs}
\cas{e)}
\begin{calculs}
u(t) &= 100t+1 &&\justif{donc $u'(t)=100$}\\[4pt]
m'(t) &= 5\times\big(-100\sin(100t+1)\big) &&\justif{la formule de $\cos(u)$}\\
m'(t) &= \res{-500\sin(100t+1)} &&
\end{calculs}
\cas{f)}
\begin{calculs}
u(x) &= 2-x &&\justif{donc $u'(x)=-1$ ; ici $n=6$}\\[4pt]
p'(x) &= 6\times u'(x)\times u(x)^{5} &&\\
p'(x) &= 6\times(-1)\times(2-x)^{5} &&\\
p'(x) &= \res{-6(2-x)^{5}} &&
\end{calculs}
\rem{En f), $u'=-1$ : le signe $-$ ne s'oublie pas, même si l'on n'écrit pas le $1$.}
\end{exo}

\begin{exo}[Un piston]
\cas{a)}
\begin{calculs}
u(t) &= 2t &&\justif{donc $u'(t)=2$}\\[4pt]
p'(t) &= 6\times\big(-2\sin(2t)\big) &&\justif{la formule de $\cos(u)$ ; le $6$ reste devant}\\
p'(t) &= \res{-12\sin(2t)} &&
\end{calculs}
\cas{b)}
\begin{calculs}
p'(0) &= -12\sin(2\times 0) &&\justif{on remplace $t$ par $0$}\\
p'(0) &= -12\times 0 &&\justif{$\sin 0=0$}\\
p'(0) &= 0 &&
\end{calculs}
\explique{À l'instant $t=0$, la vitesse du piston est $\res{0\text{ cm/s}}$ : il est
immobile.}{il est au bout de sa course, à $6$~cm}
\cas{c)}
\explique{$\sin(2t)$ est compris entre $-1$ et $1$, donc $p'(t)=-12\sin(2t)$ est compris
entre $-12$ et $12$.}{on multiplie par $-12$}
\explique{La plus grande vitesse possible du piston est $\res{12\text{ cm/s}}$ (en valeur
absolue).}{atteinte quand $\sin(2t)$ vaut $1$ ou $-1$}
\end{exo}

\begin{exo}[Dérivées en vrac]
\explique{On reconnaît d'abord la forme (somme, produit, puissance, inverse, quotient),
puis on applique la formule correspondante.}{les tableaux des propriétés 3 à 7}
\cas{a)}
\explique{Une somme : on dérive terme par terme.}{}\nopagebreak
\begin{calculs}
f'(x) &= 3\times 4x^{3}-2\times\left(-\frac{1}{x^{2}}\right) &&\justif{$\frac{1}{x}$ devient $-\frac{1}{x^{2}}$}\\
f'(x) &= \res{12x^{3}+\frac{2}{x^{2}}} &&
\end{calculs}
\cas{b)}
\explique{Un produit $u\times v$.}{}\nopagebreak
\begin{calculs}
u(x) &= x^{2}+1 &&\justif{donc $u'(x)=2x$}\\
v(x) &= x-3 &&\justif{donc $v'(x)=1$}\\[4pt]
g'(x) &= u'(x)\,v(x)+u(x)\,v'(x) &&\\
g'(x) &= 2x(x-3)+(x^{2}+1)\times 1 &&\\
g'(x) &= 2x^{2}-6x+x^{2}+1 &&\\
g'(x) &= \res{3x^{2}-6x+1} &&
\end{calculs}
\cas{c)}
\explique{Une puissance $u^{3}$.}{}\nopagebreak
\begin{calculs}
u(x) &= 2x+5 &&\justif{donc $u'(x)=2$ ; ici $n=3$}\\[4pt]
h'(x) &= 3\times u'(x)\times u(x)^{2} &&\\
h'(x) &= 3\times 2\times(2x+5)^{2} &&\\
h'(x) &= \res{6(2x+5)^{2}} &&
\end{calculs}
\cas{d)}
\explique{Un inverse : $k(x)=4\times\dfrac{1}{x^{2}+3}$.}{le $4$ reste devant}\nopagebreak
\begin{calculs}
v(x) &= x^{2}+3 &&\justif{donc $v'(x)=2x$}\\[4pt]
k'(x) &= 4\times\left(-\frac{v'(x)}{v(x)^{2}}\right) &&\\
k'(x) &= 4\times\left(-\frac{2x}{(x^{2}+3)^{2}}\right) &&\\
k'(x) &= \res{-\frac{8x}{(x^{2}+3)^{2}}} &&
\end{calculs}
\cas{e)}
\explique{Un quotient $\dfrac{u}{v}$.}{}\nopagebreak
\begin{calculs}
u(x) &= x-1 &&\justif{donc $u'(x)=1$}\\
v(x) &= x^{2}+1 &&\justif{donc $v'(x)=2x$}\\[4pt]
m'(x) &= \frac{u'(x)\,v(x)-u(x)\,v'(x)}{v(x)^{2}} &&\\
m'(x) &= \frac{1\times(x^{2}+1)-(x-1)\times 2x}{(x^{2}+1)^{2}} &&\\
m'(x) &= \frac{x^{2}+1-2x^{2}+2x}{(x^{2}+1)^{2}} &&\justif{le $-$ porte sur $2x^{2}-2x$}\\
m'(x) &= \res{\frac{-x^{2}+2x+1}{(x^{2}+1)^{2}}} &&
\end{calculs}
\cas{f)}
\explique{Un produit $u\times v$.}{}\nopagebreak
\begin{calculs}
u(x) &= x &&\justif{donc $u'(x)=1$}\\
v(x) &= \cos x &&\justif{donc $v'(x)=-\sin x$}\\[4pt]
p'(x) &= u'(x)\,v(x)+u(x)\,v'(x) &&\\
p'(x) &= 1\times\cos x+x\times(-\sin x) &&\\
p'(x) &= \res{\cos x-x\sin x} &&
\end{calculs}
\end{exo}

\partie{Dérivée et sens de variation}

\begin{exo}[Du signe de $f'$ aux variations]
\cas{Résolution}
\explique{Là où $f'(x)<0$, $f$ est décroissante ; là où $f'(x)>0$, $f$ est
croissante.}{propriété 8}
\explique{On place les images de l'énoncé au bout des flèches : $f$ décroît de $2$ à
$-3$, croît de $-3$ à $4$, puis décroît de $4$ à $1$.}{une flèche va d'une valeur à la
suivante}\nopagebreak
\begin{center}
\begin{tikzpicture}[scale=0.9]
  \tkzTabInit[lgt=2,espcl=2.2]{$x$/0.8,$f'(x)$/0.8,$f(x)$/1.6}{$-4$,$-1$,$3$,$5$}
  \tkzTabLine{,-,z,+,z,-,}
  \tkzTabVar{+/$2$,-/$-3$,+/$4$,-/$1$}
\end{tikzpicture}
\end{center}
\end{exo}

\begin{exo}[Un trinôme]
\cas{a)}
\begin{calculs}
f'(x) &= 2x-6+0 &&\justif{terme par terme}\\
f'(x) &= \res{2x-6} &&
\end{calculs}
\explique{On cherche quand $f'(x)$ s'annule :}{}\nopagebreak
\begin{calculs}
2x-6 &= 0 &&\\
2x &= 6 &&\justif{on ajoute $6$}\\
x &= 3 &&\justif{on divise par $2$}
\end{calculs}
\explique{$2x-6$ est affine, de coefficient directeur $2>0$ : négatif avant $3$, positif
après.}{signe d'une expression affine}
\cas{b)}
\explique{On calcule les images des bornes et de $3$ :}{elles complètent le tableau}\nopagebreak
\begin{calculs}
f(0) &= 5 &&\\
f(3) &= 9-18+5 = -4 &&\\
f(6) &= 36-36+5 = 5 &&
\end{calculs}
\begin{center}
\begin{tikzpicture}[scale=0.9]
  \tkzTabInit[lgt=2,espcl=2.6]{$x$/0.8,$f'(x)$/0.8,$f(x)$/1.6}{$0$,$3$,$6$}
  \tkzTabLine{,-,z,+,}
  \tkzTabVar{+/$5$,-/$-4$,+/$5$}
\end{tikzpicture}
\end{center}
\end{exo}

\begin{exo}[Un polynôme de degré 3]
\cas{a)}
\begin{calculs}
f'(x) &= 3x^{2}-3\times 2x+0 &&\justif{terme par terme}\\
f'(x) &= \res{3x^{2}-6x} &&
\end{calculs}
\cas{b)}
\explique{$f'(x)$ est un trinôme : $a=3$, $b=-6$ et $c=0$.}{on cherche ses racines
(chapitre 3)}\nopagebreak
\begin{calculs}
\Delta &= b^{2}-4ac &&\justif{le discriminant}\\
\Delta &= (-6)^{2}-4\times 3\times 0 &&\\
\Delta &= 36 &&\justif{$\Delta>0$ : deux racines}\\[4pt]
x_{1} &= \frac{6-6}{6} = 0 &&\justif{$\frac{-b-\sqrt{\Delta}}{2a}$, avec $\sqrt{36}=6$}\\
x_{2} &= \frac{6+6}{6} = 2 &&\justif{$\frac{-b+\sqrt{\Delta}}{2a}$}
\end{calculs}
\explique{$f'(x)$ est du signe de $a=3$, donc positif, à l'extérieur des racines, et
négatif entre elles.}{signe d'un trinôme}
\cas{c)}
\explique{On calcule les images des bornes et des racines :}{}\nopagebreak
\begin{calculs}
f(-1) &= -1-3+1 = -3 &&\justif{$(-1)^{3}=-1$ et $(-1)^{2}=1$}\\
f(0) &= 1 &&\\
f(2) &= 8-12+1 = -3 &&\\
f(3) &= 27-27+1 = 1 &&
\end{calculs}
\begin{center}
\begin{tikzpicture}[scale=0.9]
  \tkzTabInit[lgt=2,espcl=2.2]{$x$/0.8,$f'(x)$/0.8,$f(x)$/1.6}{$-1$,$0$,$2$,$3$}
  \tkzTabLine{,+,z,-,z,+,}
  \tkzTabVar{-/$-3$,+/$1$,-/$-3$,+/$1$}
\end{tikzpicture}
\end{center}
\rem{Plus rapide, quand $c=0$ : $3x^{2}-6x=3x(x-2)$, qui s'annule en $0$ et en $2$.}
\end{exo}

\begin{exo}[Un coefficient négatif]
\cas{a)}
\begin{calculs}
g'(x) &= -2\times 3x^{2}+3\times 2x+12-0 &&\justif{terme par terme}\\
g'(x) &= \res{-6x^{2}+6x+12} &&
\end{calculs}
\cas{b)}
\explique{$g'(x)$ est un trinôme : $a=-6$, $b=6$ et $c=12$.}{on cherche ses racines}\nopagebreak
\begin{calculs}
\Delta &= 6^{2}-4\times(-6)\times 12 &&\justif{$b^{2}-4ac$}\\
\Delta &= 36+288 &&\\
\Delta &= 324 &&\justif{$\sqrt{324}=18$}\\[4pt]
x_{1} &= \frac{-6+18}{-12} = -1 &&\justif{$2a=-12$}\\
x_{2} &= \frac{-6-18}{-12} = 2 &&
\end{calculs}
\explique{$g'(x)$ est du signe de $a=-6$, donc négatif, à l'extérieur des racines, et
positif entre elles.}{signe d'un trinôme}
\cas{c)}
\begin{calculs}
g(-2) &= 16+12-24-5 = -1 &&\justif{$-2\times(-2)^{3}=-2\times(-8)=16$}\\
g(-1) &= 2+3-12-5 = -12 &&\\
g(2) &= -16+12+24-5 = 15 &&\\
g(3) &= -54+27+36-5 = 4 &&
\end{calculs}
\begin{center}
\begin{tikzpicture}[scale=0.9]
  \tkzTabInit[lgt=2,espcl=2.2]{$x$/0.8,$g'(x)$/0.8,$g(x)$/1.6}{$-2$,$-1$,$2$,$3$}
  \tkzTabLine{,-,z,+,z,-,}
  \tkzTabVar{+/$-1$,-/$-12$,+/$15$,-/$4$}
\end{tikzpicture}
\end{center}
\rem{Avec $a<0$, le trinôme est négatif à l'extérieur des racines : les signes sont
inversés par rapport à l'exercice 42.}
\end{exo}

\begin{exo}[Un quotient]
\cas{a)}
\begin{calculs}
u(x) &= 2x+1 &&\justif{donc $u'(x)=2$}\\
v(x) &= x+1 &&\justif{donc $v'(x)=1$}\\[4pt]
f'(x) &= \frac{u'(x)\,v(x)-u(x)\,v'(x)}{v(x)^{2}} &&\justif{la formule du quotient}\\
f'(x) &= \frac{2(x+1)-(2x+1)\times 1}{(x+1)^{2}} &&\\
f'(x) &= \frac{2x+2-2x-1}{(x+1)^{2}} &&\\
f'(x) &= \frac{1}{(x+1)^{2}} &&\justif{c'est l'expression demandée}
\end{calculs}
\cas{b)}
\explique{Le numérateur $1$ est positif et le dénominateur $(x+1)^{2}$ est un carré, non
nul sur $\intff{0}{4}$ : $f'(x)>0$. Donc $f$ est \res{croissante} sur
$\intff{0}{4}$.}{propriété 8}\nopagebreak
\begin{calculs}
f(0) &= \frac{1}{1} = 1 &&\justif{images des bornes}\\
f(4) &= \frac{9}{5} &&\justif{$\frac{2\times 4+1}{4+1}$}
\end{calculs}
\begin{center}
\begin{tikzpicture}[scale=0.9]
  \tkzTabInit[lgt=2,espcl=3]{$x$/0.8,$f'(x)$/0.8,$f(x)$/1.6}{$0$,$4$}
  \tkzTabLine{,+,}
  \tkzTabVar{-/$1$,+/$\frac{9}{5}$}
\end{tikzpicture}
\end{center}
\end{exo}

\begin{exo}[Étude complète et tracé]
\cas{a)}
\begin{calculs}
f'(x) &= \res{3x^{2}-3} &&\justif{$-3x$ devient $-3$}
\end{calculs}
\explique{$f'(x)$ est un trinôme : $a=3$, $b=0$ et $c=-3$.}{}\nopagebreak
\begin{calculs}
\Delta &= 0^{2}-4\times 3\times(-3) &&\\
\Delta &= 36 &&\\[4pt]
x_{1} &= \frac{0-6}{6} = -1 &&\justif{$\sqrt{36}=6$}\\
x_{2} &= \frac{0+6}{6} = 1 &&
\end{calculs}
\explique{$f'(x)$ est du signe de $a=3$, donc positif, à l'extérieur des racines, et
négatif entre elles.}{}
\cas{b)}
\begin{calculs}
f(-2) &= -8+6 = -2 &&\\
f(-1) &= -1+3 = 2 &&\\
f(1) &= 1-3 = -2 &&\\
f(2) &= 8-6 = 2 &&
\end{calculs}
\begin{center}
\begin{tikzpicture}[scale=0.9]
  \tkzTabInit[lgt=2,espcl=2.2]{$x$/0.8,$f'(x)$/0.8,$f(x)$/1.6}{$-2$,$-1$,$1$,$2$}
  \tkzTabLine{,+,z,-,z,+,}
  \tkzTabVar{-/$-2$,+/$2$,-/$-2$,+/$2$}
\end{tikzpicture}
\end{center}
\cas{c)}
\begin{calculs}
f(0) &= 0 &&\\
f'(0) &= -3 &&\justif{$3\times 0^{2}-3$}\\[4pt]
y &= f'(0)\,(x-0)+f(0) &&\justif{la formule, avec $a=0$}\\
y &= -3x &&
\end{calculs}
\explique{La tangente $T$ a pour équation $\res{y=-3x}$.}{}
\cas{d)}
\explique{Les valeurs viennent des questions b) et c) :}{}\nopagebreak
\begin{center}
\renewcommand{\arraystretch}{1.4}
\begin{tabular}{|c|c|c|c|c|c|}
\hline
$x$ & $-2$ & $-1$ & $0$ & $1$ & $2$\\
\hline
$f(x)$ & $\res{-2}$ & $\res{2}$ & $\res{0}$ & $\res{-2}$ & $\res{2}$\\
\hline
\end{tabular}
\end{center}
\cas{e)}
\begin{minipage}[c]{0.56\linewidth}\raggedright
\explique{On trace $T$ avec deux points : l'origine, puis $(1\,;-3)$ (on avance de $1$,
on descend de $3$).}{$y=-3x$}
\explique{On place les cinq points du tableau de valeurs, puis on trace $\Cf$ : tangente
horizontale en $-1$ et en $1$ (extremums), tangente $T$ en $0$.}{la courbe suit le
tableau de variations}
\end{minipage}\hfill
\begin{minipage}[c]{0.4\linewidth}\centering
\begin{tikzpicture}[x=0.6cm,y=0.6cm]
  \repere{-3}{-4}{3}{4}
  \gradx{-2,-1,1,2} \grady{-3,-2,-1,1,2,3}
  \draw[coulCourbeB,thick] (-1.3,3.9)--(1.3,-3.9);
  \draw[coulCourbe,line width=1.1pt]
    (-2,-2) .. controls (-1.667,1) and (-1.333,2) .. (-1,2)
            .. controls (-0.667,2) and (-0.333,1) .. (0,0)
            .. controls (0.333,-1) and (0.667,-2) .. (1,-2)
            .. controls (1.333,-2) and (1.667,-1) .. (2,2);
  \path[coulCourbe] (-2,-2) node[croix]{} (-1,2) node[croix]{} (0,0) node[croix]{}
    (1,-2) node[croix]{} (2,2) node[croix]{};
  \path[coulCourbeB] (1,-3) node[croix]{};
  \node[coulCourbeB,font=\footnotesize] at (-1.5,3.3) {$T$};
  \node[coulCourbe,font=\footnotesize] at (2.4,1.2) {$\Cf$};
\end{tikzpicture}
\end{minipage}
\end{exo}

\partie{Problèmes concrets : maximum et minimum}

\begin{exo}[Un coût minimal]
\cas{a)}
\begin{calculs}
C'(q) &= 2q-30+0 &&\justif{la variable s'appelle $q$ au lieu de $x$}\\
C'(q) &= \res{2q-30} &&
\end{calculs}
\cas{b)}
\explique{On cherche quand $C'(q)$ s'annule :}{}\nopagebreak
\begin{calculs}
2q-30 &= 0 &&\\
2q &= 30 &&\\
q &= 15 &&
\end{calculs}
\explique{$2q-30$ est affine, de coefficient directeur $2>0$ : négatif avant $15$,
positif après.}{signe d'une expression affine}\nopagebreak
\begin{calculs}
C(0) &= 300 &&\\
C(15) &= 225-450+300 = 75 &&\justif{$15^{2}-30\times 15+300$}\\
C(25) &= 625-750+300 = 175 &&
\end{calculs}
\begin{center}
\begin{tikzpicture}[scale=0.9]
  \tkzTabInit[lgt=2,espcl=2.6]{$q$/0.8,$C'(q)$/0.8,$C(q)$/1.6}{$0$,$15$,$25$}
  \tkzTabLine{,-,z,+,}
  \tkzTabVar{+/$300$,-/$75$,+/$175$}
\end{tikzpicture}
\end{center}
\cas{c)}
\explique{Le coût est minimal pour $q=15$ dizaines, soit \res{150 pièces} ; il vaut
alors \res{75~\euro}.}{$C'$ s'annule en $15$ en passant de $-$ à $+$ : un minimum}
\rem{$q$ compte des dizaines de pièces : on répond à la question posée, en pièces.}
\end{exo}

\begin{exo}[La boîte la plus grande]
\cas{a)}
\begin{calculs}
V(x) &= x(12-2x)^{2} &&\justif{on recopie l'expression}\\
V(x) &= x(144-48x+4x^{2}) &&\justif{$(a-b)^{2}=a^{2}-2ab+b^{2}$ : $2\times 12\times 2x=48x$}\\
V(x) &= 144x-48x^{2}+4x^{3} &&\justif{on distribue $x$}\\
V(x) &= 4x^{3}-48x^{2}+144x &&\justif{c'est l'expression demandée}
\end{calculs}
\cas{b)}
\begin{calculs}
V'(x) &= 4\times 3x^{2}-48\times 2x+144 &&\justif{terme par terme}\\
V'(x) &= \res{12x^{2}-96x+144} &&
\end{calculs}
\cas{c)}
\explique{$V'(x)$ est un trinôme : $a=12$, $b=-96$ et $c=144$.}{}\nopagebreak
\begin{calculs}
\Delta &= (-96)^{2}-4\times 12\times 144 &&\\
\Delta &= 9\,216-6\,912 &&\\
\Delta &= 2\,304 &&\justif{$\sqrt{2\,304}=48$}\\[4pt]
x_{1} &= \frac{96-48}{24} = 2 &&\\
x_{2} &= \frac{96+48}{24} = 6 &&
\end{calculs}
\explique{$V'(x)$ est du signe de $a=12$, donc positif, à l'extérieur des racines, et
négatif entre elles : positif sur $\intfo{0}{2}$, négatif sur $\intoo{2}{6}$.}{$6$ est
à la fois une racine et la borne}
\rem{Plus simple : $V'(x)=12(x^{2}-8x+12)$, et l'on cherche les racines de
$x^{2}-8x+12$, avec des nombres plus petits.}
\cas{d)}
\begin{calculs}
V(0) &= 0 &&\\
V(2) &= 2\times 8^{2} = 128 &&\justif{avec la forme de départ : $12-2\times 2=8$}\\
V(6) &= 6\times 0^{2} = 0 &&
\end{calculs}
\begin{center}
\begin{tikzpicture}[scale=0.9]
  \tkzTabInit[lgt=2,espcl=2.6]{$x$/0.8,$V'(x)$/0.8,$V(x)$/1.6}{$0$,$2$,$6$}
  \tkzTabLine{,+,z,-,z}
  \tkzTabVar{-/$0$,+/$128$,-/$0$}
\end{tikzpicture}
\end{center}
\cas{e)}
\explique{Le volume est maximal pour $\res{x=2}$~cm ; il vaut alors
$\res{128\text{ cm}^{3}}$.}{$V'$ s'annule en $2$ en passant de $+$ à $-$ : un maximum}
\end{exo}

\begin{exo}[Le tambour d'une machine à laver]
\cas{a)}
\begin{calculs}
N'(t) &= -0{,}5\times 2t+20 &&\justif{la variable s'appelle $t$}\\
N'(t) &= \res{-t+20} &&
\end{calculs}
\cas{b)}
\explique{Le tambour s'arrête quand sa vitesse de rotation est nulle :}{$N'(t)=0$}\nopagebreak
\begin{calculs}
-t+20 &= 0 &&\\
-t &= -20 &&\\
t &= 20 &&
\end{calculs}
\explique{Le tambour s'arrête au bout de $\res{20}$ secondes.}{avant, $-t+20>0$ : il
tourne}
\cas{c)}
\begin{calculs}
N(20) &= -0{,}5\times 20^{2}+20\times 20 &&\justif{on remplace $t$ par $20$}\\
N(20) &= -200+400 &&\\
N(20) &= 200 &&
\end{calculs}
\explique{Jusqu'à l'arrêt, le tambour a fait $\res{200}$ tours.}{$N$ est maximal en
$20$ : le nombre de tours n'augmente plus}
\end{exo}

\partie{Résoudre une équation $f(x)=k$}

\begin{exo}[Nombre de solutions]
\explique{Sur chaque intervalle où $f$ varie dans un seul sens, l'équation $f(x)=k$ a une
solution si $k$ est compris entre les deux valeurs aux bouts de la flèche, aucune
sinon.}{propriété 10, intervalle par intervalle}
\cas{a)}
\explique{$0$ est entre $-2$ et $5$, entre $5$ et $-4$, entre $-4$ et $1$ : une solution
sur chacun des trois intervalles. $\res{3}$ solutions, dans $\intoo{-3}{0}$,
$\intoo{0}{2}$ et $\intoo{2}{4}$.}{trois flèches qui croisent la valeur $0$}
\cas{b)}
\explique{$3$ est entre $-2$ et $5$, entre $5$ et $-4$, mais pas entre $-4$ et $1$.
$\res{2}$ solutions, dans $\intoo{-3}{0}$ et $\intoo{0}{2}$.}{sur $\intff{2}{4}$, $f$ ne
dépasse pas $1$}
\cas{c)}
\explique{$-3$ n'est pas entre $-2$ et $5$, mais il est entre $5$ et $-4$ et entre $-4$ et
$1$. $\res{2}$ solutions, dans $\intoo{0}{2}$ et $\intoo{2}{4}$.}{sur $\intff{-3}{0}$, $f$
reste au-dessus de $-2$}
\cas{d)}
\explique{Le maximum de $f$ est $5$ : $f(x)$ n'atteint jamais $6$. \res{aucune} solution.}{$6$
n'est entre les bouts d'aucune flèche}
\end{exo}

\begin{exo}[Une seule solution]
\cas{a)}
\begin{calculs}
f'(x) &= 3x^{2}+1 &&\justif{terme par terme}
\end{calculs}
\explique{Un carré est positif, donc $3x^{2}+1>0$ : $f$ est strictement croissante sur
$\intff{0}{2}$.}{$f'(x)>0$ : $f$ croît}
\cas{b)}
\begin{calculs}
f(0) &= -3 &&\justif{les images des bornes}\\
f(2) &= 8+2-3 = 7 &&
\end{calculs}
\explique{$f$ est strictement croissante et $0$ est compris entre $-3$ et $7$ :
l'équation $f(x)=0$ a une seule solution $\alpha$ dans $\intff{0}{2}$.}{théorème des
valeurs intermédiaires}
\cas{c)}
\explique{Tableau de valeurs avec un pas de $1$ : $f(1)=-1$ et $f(2)=7$. Donc
$1<\alpha<2$.}{$f(x)$ passe de négatif à positif}
\explique{Pas de $0{,}1$ à partir de $1$ : $f(1{,}2)=-0{,}072$ et $f(1{,}3)=0{,}497$.
Donc $\res{1{,}2<\alpha<1{,}3}$.}{on recommence entre $1$ et $2$}
\end{exo}

\begin{exo}[Trois solutions]
\cas{a)}
\begin{calculs}
g'(x) &= \res{3x^{2}-12} &&\justif{$-12x$ devient $-12$}
\end{calculs}
\explique{$g'(x)$ est un trinôme : $a=3$, $b=0$ et $c=-12$.}{}\nopagebreak
\begin{calculs}
\Delta &= 0^{2}-4\times 3\times(-12) &&\\
\Delta &= 144 &&\justif{$\sqrt{144}=12$}\\[4pt]
x_{1} &= \frac{0-12}{6} = -2 &&\\
x_{2} &= \frac{0+12}{6} = 2 &&
\end{calculs}
\explique{$g'(x)$ est du signe de $a=3$, donc positif, à l'extérieur des racines, et
négatif entre elles.}{signe d'un trinôme}
\cas{b)}
\begin{calculs}
g(-4) &= -64+48 = -16 &&\\
g(-2) &= -8+24 = 16 &&\\
g(2) &= 8-24 = -16 &&\\
g(4) &= 64-48 = 16 &&
\end{calculs}
\begin{center}
\begin{tikzpicture}[scale=0.9]
  \tkzTabInit[lgt=2,espcl=2.2]{$x$/0.8,$g'(x)$/0.8,$g(x)$/1.6}{$-4$,$-2$,$2$,$4$}
  \tkzTabLine{,+,z,-,z,+,}
  \tkzTabVar{-/$-16$,+/$16$,-/$-16$,+/$16$}
\end{tikzpicture}
\end{center}
\cas{c)}
\explique{Sur chacun des trois intervalles, $g$ varie dans un seul sens entre $-16$ et
$16$, et $5$ est compris entre $-16$ et $16$ : une solution sur chacun. L'équation
$g(x)=5$ a $\res{3}$ solutions dans $\intff{-4}{4}$.}{propriété 10, trois fois}
\cas{d)}
\explique{La plus grande solution, notée $\alpha$, est dans $\intff{2}{4}$, où $g$ est
strictement croissante.}{}
\explique{Pas de $1$ : $g(3)=-9$ et $g(4)=16$. Donc $3<\alpha<4$.}{$g(x)$ passe
au-dessus de $5$}
\explique{Pas de $0{,}1$ à partir de $3$ : $g(3{,}6)=3{,}456$ et $g(3{,}7)=6{,}253$.
Donc $\res{3{,}6<\alpha<3{,}7}$.}{on compare à $5$, et non à $0$}
\rem{Erreur fréquente : chercher où $g(x)$ change de signe, comme pour $g(x)=0$. Ici, on
cherche où $g(x)$ dépasse $5$.}
\end{exo}

\end{document}
